\documentclass[11pt,letterpaper]{article}
\usepackage[T1]{fontenc}
\usepackage[utf8]{inputenc}
\usepackage{lmodern}
\usepackage[margin=1in,headheight=15pt]{geometry}
\usepackage{amsmath,amssymb,amsthm,mathtools,mathrsfs}
\usepackage{microtype,booktabs,longtable,tabularx,array,enumitem}
\usepackage[dvipsnames]{xcolor}
\usepackage{listings,fancyhdr}
\usepackage[colorlinks=true,linkcolor=MidnightBlue,citecolor=MidnightBlue,urlcolor=MidnightBlue]{hyperref}
\usepackage[nameinlink,noabbrev]{cleveref}
\hypersetup{pdftitle={Surreal and Surcomplex Numbers in Symbolic Computer Algebra},pdfauthor={Research and software-design exposition},pdfsubject={Exact representations, effective subfields, Hahn supports, certified precision, and Wolfram Language}}
\pagestyle{fancy}\fancyhf{}
\fancyhead[L]{\small Surreal and surcomplex computer algebra}
\fancyhead[R]{\small Representation, algorithms, and semantics}
\fancyfoot[C]{\thepage}
\setlength{\parskip}{0.28em}
\setlength{\parindent}{1.2em}
\setlist{leftmargin=2em,itemsep=0.15em,topsep=0.3em}
\setcounter{tocdepth}{2}
\numberwithin{equation}{section}
\newtheorem{theorem}{Theorem}[section]
\newtheorem{proposition}[theorem]{Proposition}
\newtheorem{lemma}[theorem]{Lemma}
\newtheorem{corollary}[theorem]{Corollary}
\theoremstyle{definition}
\newtheorem{definition}[theorem]{Definition}
\newtheorem{example}[theorem]{Example}
\theoremstyle{remark}\newtheorem{remark}[theorem]{Remark}
\newcommand{\No}{\mathbf{No}}
\newcommand{\SC}{\mathbf{SC}}
\newcommand{\R}{\mathbb R}
\newcommand{\C}{\mathbb C}
\newcommand{\Q}{\mathbb Q}
\newcommand{\N}{\mathbb N}
\newcommand{\Z}{\mathbb Z}
\newcommand{\RA}{\mathbb R_{\rm alg}}
\newcommand{\CA}{\overline{\mathbb Q}}
\newcommand{\OO}{\mathcal O}
\newcommand{\mm}{\mathfrak m}
\newcommand{\HH}{\mathscr H}
\newcommand{\st}{\operatorname{st}}
\newcommand{\supp}{\operatorname{supp}}
\newcommand{\lc}{\operatorname{lc}}
\newcommand{\ord}{\operatorname{ord}}
\newcommand{\Res}{\operatorname{Res}}
\newcommand{\Tr}{\operatorname{Tr}}
\newcommand{\Disc}{\operatorname{Disc}}
\newcommand{\sh}{^{\sharp}}
\newcommand{\Oh}[1]{O_v(t^{#1})}
\DeclareMathOperator{\Arg}{Arg}
\DeclareMathOperator{\Log}{Log}
\newcolumntype{P}[1]{>{\raggedright\arraybackslash}p{#1}}
\newcolumntype{Y}{>{\raggedright\arraybackslash}X}
\lstdefinestyle{code}{basicstyle=\ttfamily\small,columns=fullflexible,keepspaces=true,breaklines=true,showstringspaces=false,frame=single,framerule=0.25pt,rulecolor=\color{gray},xleftmargin=0.3em,xrightmargin=0.3em,aboveskip=0.7em,belowskip=0.7em}
\lstset{style=code}
\begin{document}
\hypersetup{pageanchor=false}
\begin{titlepage}
\centering
\vspace*{0.7cm}
{\large\scshape Mathematical foundations and an implementation design\par}
\vspace{1cm}
{\Huge\bfseries Surreal and Surcomplex Numbers\par}
\vspace{0.45cm}
{\LARGE in Symbolic Computer Algebra\par}
\vspace{0.65cm}
{\Large Exact representations, effective subfields,\par certified series, and Wolfram Language\par}
\vspace{0.85cm}
{\large September 21, 2026\par}
\vspace{0.8cm}
\begin{minipage}{0.93\textwidth}
\begin{abstract}
Surreal numbers can be supported meaningfully in a symbolic computer algebra
system without pretending to store every surreal or to decide every identity.
The appropriate design combines several exact representations: sparse Conway
normal forms, rational functions in ordered monomials, algebraic root
specifications, certified generalized-series generators, ordinary analytic
lifts, and selected transserial expressions. Each representation has a different
algorithmic contract.

We identify a countable, effectively presented real-closed subfield containing
arbitrarily many finite levels of infinitesimal scales; its complexification
supports exact polynomial algebra and algebraic root objects. We prove explicit
limitations for arbitrary program-described coefficients and support validation,
and explain why a finite Hahn coefficient calculation need not provide a finite
truncation at a requested valuation. A constructive positive-grid lemma separates
these two issues. We develop precision propagation, root-cluster algorithms,
local analytic evaluation, implicit functions, and collision-stable residues.

The three supplied manuscripts determine the meaning of Hahn coherence,
strong summability, and the distinction between fine-local and all-scale
functions. The proposed software architecture preserves those distinctions,
including the nonuniqueness of a global surcomplex exponential at infinite
imaginary arguments. Wolfram Language facilities are assessed against current
official documentation. A small, separately delimited Wolfram Language prototype
and reproducible tests demonstrate exact sparse arithmetic; they do not purport
to implement the entire architecture.
\end{abstract}
\end{minipage}
\vfill
\begin{minipage}{0.93\textwidth}\small
\textbf{Status.} This is an implementation-oriented mathematical exposition.
The supplied research manuscripts are treated as sources, not as independently
refereed or machine-verified publications. Established theorems, deductions
proved here, proposed interfaces, and executed example tests are distinguished.
No universal surreal simplifier or complete surcomplex analytic system is claimed.
\end{minipage}
\end{titlepage}
\hypersetup{pageanchor=true}
\pagenumbering{roman}
\tableofcontents
\clearpage\pagenumbering{arabic}

\section{The answer: a layered exact system, not one universal number type}
\label{sec:answer}

\subsection{What should be supported first?}
A useful implementation should begin with an exact ordered field of rational
expressions in explicitly ordered surreal monomials, extend it by algebraic
root objects, and attach certified series representations when expansions are
wanted. This already permits genuine infinite and infinitesimal quantities,
not merely a symbolic placeholder for a limit. Complexifying the real field
gives an exact surcomplex coefficient domain. Polynomial division, greatest
common divisors, resultants, finite linear algebra, and polynomial root
specifications then become natural operations.

A second layer should represent exact infinite series by finite rules with
proved support conditions. It can evaluate Taylor series at infinitesimals,
compute requested coefficients, and, where justified, return finite
truncations with valuation error bounds. A third layer can add selected
exponential--logarithmic and transserial constructors, and ordinary holomorphic
functions lifted to finite surcomplex arguments. These are extensions with
additional semantics, not consequences of adding one arithmetic head.

The design proposed here has a deliberately asymmetric policy: exact algebra
is aggressive inside a declared decidable domain, while more general analytic
expressions remain symbolic unless their domain, support, and branch obligations
have been discharged. An unevaluated but well-defined expression is a legitimate
exact result. An expression with an unproved summability obligation is not yet
an admitted exact number.

\subsection{Four different meanings of ``exact''}
\label{subsec:exact}
\begin{longtable}{@{}P{0.21\textwidth}P{0.35\textwidth}P{0.36\textwidth}@{}}
\toprule
Property & What it promises & What it does not promise\\
\midrule\endhead
Exact denotation & A finite description singles out one mathematical value. &
A normal form, a numerical approximation, or a terminating equality test.\\[0.5em]
Effective coefficient access & A requested coefficient can be calculated by a
terminating procedure. & The first nonzero coefficient can always be found,
or all coefficients below a cutoff can be listed.\\[0.5em]
Decidable algebra & Equality and the required field operations terminate on
all valid inputs in a declared class. & Closure under every elementary or
special function.\\[0.5em]
Certified approximation & The answer includes a proved error specification
in a named topology or valuation. & Equality to the truncated expression, or
approximation in the full surreal fine topology.\\
\bottomrule
\end{longtable}
For example, $(1-t)^{-1}$ is a finite exact field expression. Its Hahn normal
form has infinitely many terms. A displayed polynomial
$1+t+\cdots+t^9$ is not that same exact number; it is an approximation with a
specific remainder. Conversely, the program description of a series can be
exact while equality of two such programs' values is undecidable. These are
structural distinctions, not implementation inconveniences.

\subsection{How the supplied manuscripts are used}
\label{subsec:sources}
The supplied \emph{Surcomplex Polynomial Algebra} \cite{SourcePolynomial}
provides the set-sized Hahn workspaces, initial-polynomial root counts,
Newton profiles, support-controlled coprime lifting, perturbation estimates,
and finite residue algebras used below. It explicitly warns that its finite
algebraic constructions are not automatically effective for arbitrary Hahn
inputs. We turn that warning into concrete representation contracts.

The supplied \emph{Hartogs Coherence, Finite Maps, and Residue Duality}
\cite{SourceResearch} provides the fixed-domain coefficient rings,
positive-operator inverses, coupled finite-free systems, and support bounds
$S+E^*$. Its mathematical common-support and division theorems do not, by
themselves, synthesize finite certificates or decide identities of arbitrary
holomorphic coefficients. Those missing computational assumptions are made
explicit rather than attributed to the source.

The supplied \emph{Surcomplex Analysis} \cite{SourceAnalysis} distinguishes
fine-local formal germs, Hahn-coherent functions on ordinary-domain
thickenings, and all-scale coherent functions. We preserve this terminology.
In particular, its all-scale polynomial rigidity result is a statement about
its specified function category, not a statement that surreal exponentials
cannot exist. Its explicit family of global exponentials is important for
software branch and extension metadata.

The polynomial manuscript also cites a fourth, global-divisor manuscript.
That document is not among the three supplied files and is not used here as
an independently inspected source. External references in this article supply
published foundations and computational precedents. The proposed data types,
capability system, and development stages are design recommendations of the
present article, not claims that the supplied manuscripts already implement them.

\section{The mathematical objects that the software must distinguish}
\label{sec:objects}

\subsection{Set-sized workspaces and their effective counterparts}
Put $\SC=\No[i]$. Following the sources, write
\[
 t^\gamma=\omega^{-\gamma},\qquad
 F_\Gamma=\R((t^\Gamma)),\qquad K_\Gamma=\C((t^\Gamma)),
\]
where $\Gamma$ is a set-sized ordered additive subgroup of $\No$. A Hahn series
is an expression
\begin{equation}\label{eq:hahn}
 A=\sum_{\gamma\in S}a_\gamma t^\gamma
\end{equation}
with well-ordered nonzero support $S$. The Conway normal-form embedding gives
these expressions their surreal or surcomplex meanings. When $\Gamma$ is
divisible, the real field is real closed and the complex field is algebraically
closed; the Hahn-field algebraic-closure input is classical
\cite{Poonen,SourcePolynomial}.

This is a semantic workspace, not a finite data structure. Even
$\Q((t))$ contains an uncountable collection of arbitrary rational coefficient
sequences. A practical implementation replaces $\R$ or $\C$ by a declared
coefficient representation, replaces arbitrary exponents by an effective
ordered-group representation, and admits only selected finitely described
supports and coefficient rules. Such a replacement may preserve some closure
properties and lose others. Each closure claim must refer to the replacement,
not merely to the surrounding Hahn field.

\subsection{Conway monomials are not generic powers}
\label{subsec:omega}
The symbol $\omega^a$ in a Conway normal form denotes the omega-map. It obeys
\[
 \omega^{a+b}=\omega^a\omega^b,
\]
but it is not defined by the generic real-power prescription
$\exp_{\No}(a\log_{\No}\omega)$. The canonical real surreal exponential and
logarithm constitute an additional structure \cite{Gonshor,vdDE}.

The distinction is substantial for exponents of the omega-map that are
surreal but not ordinary real numbers. For instance $\omega^{1/\omega}$, in Conway notation,
is an infinite monomial: its exponent is positive, so its Archimedean class
strictly exceeds that of $1$. In contrast,
\[
 \exp_{\No}\!\left(\frac{\log_{\No}\omega}{\omega}\right)
 =1+\text{a positive infinitesimal}.
\]
Here $\log_{\No}\omega/\omega$ is infinitesimal, as follows from the growth
properties of the surreal exponential. Rational powers of $\omega$ cause no
such conflict: both constructions give the unique appropriate positive
algebraic root. The collision of notation becomes dangerous when a package
extends that rational-power simplification to arbitrary surreal exponents.

A defensible internal syntax therefore has distinct constructors such as
\texttt{ConwayMonomial[a]} and \texttt{SurrealExp[a]}. A printed superscript is
not sufficient semantic metadata. General complex powers require still more
information: a chosen logarithm or a specified algebraic-root branch.

\subsection{Order, valuation, and standard part}
For nonzero $A$ define $v(A)=\min\supp(A)$ and
$\lc(A)=a_{v(A)}$, and set $v(0)=+\infty$. Then
\[
 v(AB)=v(A)+v(B),\qquad v(A+B)\ge\min\{v(A),v(B)\}.
\]
For a real series, its sign is the sign of its leading coefficient. For a
surcomplex number $z=x+iy$, conjugation and modulus are
\[
 \bar z=x-iy,\qquad |z|=\sqrt{x^2+y^2}.
\]
The modulus satisfies the ordinary triangle inequality, while the valuation
satisfies an ultrametric inequality. An order disk $|z-a|<r$ is not a
valuation ball $v(z-a)>v(r)$ \cite{SourcePolynomial}.

With ordinary Archimedean coefficients, $v(z)\ge0$ means that $z$ is finite,
and $v(z)>0$ means that it is infinitesimal. On finite elements,
$\st(z)$ is the coefficient at exponent zero. The zero coefficient of an
infinite element can also be extracted, but calling it a standard part would
misstate the domain of the residue homomorphism. For coefficients in an
Archimedean subfield such as $\RA$ or $\CA$, the same assertions hold with
residue values in that coefficient field.

Software must distinguish an exact valuation from a lower bound. If the
leading displayed coefficient is not known to be nonzero, it has not yet
established the valuation. Cancellation in a sum is not a cosmetic issue:
it can change sign, scale, and the number of roots in a selected ball.

\subsection{Three incompatible uses of ``infinitesimal''}
A nonzero surreal infinitesimal has an inverse and is not nilpotent. A dual
number $a+b\epsilon$ with $\epsilon^2=0$ belongs to a different algebra.
A truncated series modulo $t^N$ also belongs to a quotient ring with
nilpotents. These rings are useful for jets and automatic differentiation,
but neither is a subfield of the surreals. In particular, storing a surreal
$t$ in a dual-number object and applying its multiplication law would assert
$t^2=0$, an incorrect equality.

An asymptotic variable tending to zero is a fourth kind of object: it is a
variable in a function or germ, not yet a particular surreal constant.
Converting its formal expansion into an exact Hahn value requires a selected
interpretation and a summability proof. We retain these distinctions even
when the same displayed expression is convenient for all four purposes.

\section{Limits of finite representation and universal simplification}
\label{sec:limits}

\subsection{There is no finite notation for every surreal}
\begin{proposition}[The finite-description bound]\label{prop:countable}
For any fixed finite or countable alphabet, a language of finite strings
can denote at most countably many individually specified numbers. It cannot
denote every real number, and therefore cannot denote every surreal number.
\end{proposition}
\begin{proof}
There are countably many strings of each finite length and countably many
lengths. Taking the image of this countable set under a denotation map
cannot increase cardinality. The reals are uncountable and are contained in
$\No$.
\end{proof}
Allowing a symbol for an arbitrary unspecified real parameter gives a symbolic
family, not a finite encoding of every possible value of that parameter.
Allowing an oracle for an arbitrary real moves information outside the finite
input. A proof assistant can quantify over all surreals without making them
all machine-representable data. Likewise, a computer algebra system can state
universal identities in abstract variables without materializing their values.

There is no uniquely largest useful finitely representable subset. Adding one
new constant, one new support notation, or one new certified function may
strictly extend an existing language. The practical question is which
representations offer a useful and honest collection of operations.

\subsection{A minimal undecidability obstruction}
\begin{theorem}[Total coefficient access does not decide zero]
\label{thm:undecidable}
There is no algorithm which, for every total program specifying rational
coefficients $a_n$ of a valid series $A=\sum_{n\ge0}a_nt^n$, decides whether
$A=0$. This remains impossible when every input in the reduction actually
denotes either zero or a single monomial.
\end{theorem}
\begin{proof}
Given a Turing machine $M$ and its input, define $a_n$ to be $1$ if $M$
first halts at step $n$, and $0$ otherwise. The value of $a_n$ is computable
by a bounded simulation, so the coefficient program is total. Its series has
empty support if $M$ never halts, and singleton support if it halts. Thus it
is always a valid Hahn series, indeed a polynomial. A zero test for its
program representation would decide whether $M$ halts.
\end{proof}
The issue is the representation, not only the mathematical subset of values.
Explicit finite sparse polynomials have decidable equality. Opaque programs
promised merely to denote finite sparse polynomials need not have a computable
support bound. They can describe exactly the same set of values with a much
weaker computational contract.

This obstruction is independent of difficult identities among elementary
constants. Richardson's work supplies other undecidability results for
sufficiently rich expression languages of real functions
\cite{Richardson}. Those results must not be inflated into a blanket claim
that every restricted exponential-constant identity problem is known to be
undecidable. The elementary coefficient-program reduction above is enough
for the unrestricted Hahn-generator interface considered here.

\subsection{Support admissibility is not a universal automatic test}
\begin{proposition}[Unrestricted support validation is undecidable]
\label{prop:validation}
No algorithm decides Hahn admissibility for every program which outputs a
candidate support using total bounded-time coefficient tests.
\end{proposition}
\begin{proof}
For a machine $M$, assign a coefficient at exponent $-n$, $n\ge1$, equal to
$1$ if $M$ has halted by step $n$, and $0$ otherwise. Every individual test
terminates. If $M$ never halts, the exact support is empty and admissible.
If it halts, the support contains an infinite strictly decreasing tail
$-n,-(n+1),\ldots$ and is not well ordered. Deciding admissibility would
therefore decide halting.
\end{proof}
The remedy is not to omit validation. It is to use constructors whose support
properties follow from their syntax or from checked certificates, and to treat
unrestricted user generators as untrusted until obligations are supplied.
A validator can be total and complete for a restricted grammar; it cannot be
both total and complete for all programs of the preceding proposition.

\section{A taxonomy of exact representations}
\label{sec:representations}

\subsection{Finite cuts and sign expansions}
The original cut presentation and the sign-sequence presentation retain
simplicity and birthday information. Fully finite cut trees describe the
finite-birthday numbers, namely the dyadic rationals. They are excellent
for teaching the recursive construction and checking small identities, but
not a suitable numerical core: even $1/3$ requires more than a fully finite
cut tree. Division therefore leaves this representation.

Compressed infinite cuts can be exact, for example the familiar description
of $\omega$ by all natural numbers on the left. Their effectiveness depends
on the language used for the options, not on the braces. Computing the
simplest number between arbitrary program-defined classes of options is not
a generic finite operation. A cut constructor should record an ordinal or
well-foundedness certificate and an inequality certificate; it should not
accept an arbitrary predicate and promise a canonical normal form.

Finite ordinal notations can describe many infinite surreals, but ordinal
addition and multiplication are not surreal field addition and multiplication.
For embedded ordinals the latter agree with natural, or Hessenberg,
arithmetic. A package must therefore distinguish \texttt{OrdinalPlus} from
surreal addition. Birthday and simplicity queries should remain a separate
capability from field equality.

\subsection{Finite normal forms and rational expressions}
A sparse normal form
\[
 A=c_1t^{\gamma_1}+\cdots+c_mt^{\gamma_m},\qquad
 \gamma_1<\cdots<\gamma_m,
\]
is a particularly good exact representation when coefficient and exponent
equality are decidable. Addition merges supports; multiplication forms pairwise
exponent sums and collects equal terms. Conjugation acts on coefficients.
For real coefficients, ordering is a leading-term calculation.

Finite support is not closed under inversion:
\[
 \frac1{1-t}=\sum_{n\ge0}t^n.
\]
It is therefore useful to preserve a numerator and denominator as a rational
expression instead of forcing expansion. If both are sparse polynomials in
independent monomial generators, equality is decided by cross multiplication.
The result is an exact field element even when no finite normal form exists.
A common-denominator normal form and polynomial gcd cancellation can improve
performance, but are not necessary for the correctness of this equality test.

\subsection{Algebraic specifications and branch selectors}
An algebraic object should contain a polynomial over the declared coefficient
field and enough data to select one root. A bare equation $P(z)=0$ denotes a
finite set or a relation, not one number. Valid selectors include an ordered
real root index, a Thom encoding, a certified simple residual root after a
scale change, or a combination of such data. A valuation ball containing a
multiple cluster is not a unique-root selector.

An exact algebraic number need not have a finite radical expression. The same
is true over surreal coefficient fields. Implicit root objects are therefore
a central representation, not an unfortunate fallback. Their advantages include
finite descriptions of infinitely supported roots and algebraic closure of
a well-chosen effective core. Their costs include root-branch management and
potentially large algebraic degrees.

\subsection{Formal-series rules, equations, and recurrences}
The description
\[
 A=\sum_{n\ge0}n!t^n
\]
is finite and exact as a Hahn value. So are descriptions using a rational
generating function, a total coefficient recurrence with specified initial
data, or a formal implicit equation with a uniqueness certificate. These
mechanisms have different equality algorithms. Rational generating functions
have rational-function equality; algebraic series have algebraic equality;
arbitrary coefficient programs do not have a universal zero test.

A useful intermediate class consists of formal solutions of specified linear
differential equations with enough initial data. For an ordinary point, the
initial jet and the recurrence determine the solution. Within a representation
that can construct an annihilating equation for a difference, identity can be
reduced to a finite initial-data test. Singular points require a certified
recurrence start and treatment of exceptional indices; the phrase
``satisfies a differential equation'' alone does not guarantee uniqueness.
Products and derivatives can be handled in the usual D-finite setting, but
reciprocals and arbitrary composition need not remain D-finite. Wolfram's
\texttt{DifferentialRoot} is a relevant ordinary-function precedent, not an
automatic Hahn-series implementation \cite{WLDifferentialRoot}.

\subsection{Grids, nested towers, and unrestricted Hahn data}
A positive grid has the form
\begin{equation}\label{eq:grid}
 a+\N e_1+\cdots+\N e_m,\qquad e_j>0.
\end{equation}
The coefficient rule on this grid is still part of the input. Finite grids
are useful because the support condition follows from positivity and the
Hahn--Neumann lemma, and coefficient extraction can be made effective under
explicit assumptions. Grid-based transseries have a substantial computational
literature; they should not be presented as a newly invented general method
\cite{vdH}.

Recursive exponent towers can represent, for example, finitely many nested
Conway monomials whose exponents themselves have finite normal forms. A
well-founded finite dependency graph can make arithmetic and comparison
recursive. Such a language captures far more than rational powers, but every
extension must specify normalization of exponents, cross-level coercions,
and which operations preserve the syntax. A general expression in
\texttt{SurrealExp}, \texttt{SurrealLog}, and \texttt{ConwayMonomial} need not
possess a known complete normalizer merely because its expression graph is finite.

At the broadest level one can store an abstract Hahn datum with a proof of
well-ordered support and coefficient specifications. This supports theorem
statements and some symbolic transformations. It should advertise only those
computational capabilities actually supplied by the certificates, rather than
pretending that a proof of existence is an enumeration algorithm.

\subsection{A representation and closure map}
\begin{longtable}{@{}P{0.24\textwidth}P{0.31\textwidth}P{0.37\textwidth}@{}}
\toprule
Representation & Exact values included & Main closure or decision boundary\\
\midrule\endhead
Finite cut trees & Dyadic rationals; finite-birthday data & Not closed under
division. Simplicity information is retained.\\[0.45em]
Finite sparse normal forms & Finite sums of admitted monomials & Closed under
addition and multiplication; inverse usually needs another representation.\\[0.45em]
Monomial rational functions & Fractions of those sums & Exact field algebra
and equality with effective independent generators.\\[0.45em]
Effective real closure and its complexification & All algebraic elements over
the chosen ordered monomial field & Real order and algebraic equality are
decidable; polynomial roots remain representable.\\[0.45em]
Rational or algebraic series & Infinite supports defined by finite algebraic
data & Strong exact structure; root selection remains necessary.\\[0.45em]
Certified recurrence or D-finite series & Many special-function expansions,
including some divergent formal series & Coefficients effective with adequate
initial data; closure and identity are representation-specific.\\[0.45em]
Certified positive grids & Multiscale series with finite generators & Fixed
coefficients can be effective; a bounded valuation prefix can still be infinite.\\[0.45em]
Selected transserial syntax & Exponential and logarithmic scale hierarchies &
Requires a specified ordered differential and composition framework.\\[0.45em]
General program or proof-described Hahn data & A broad countable collection
of exact values & No universal equality, support-validation, or truncation
algorithm.\\
\bottomrule
\end{longtable}

\section{An effective real-closed and algebraically closed core}
\label{sec:core}

\subsection{A concrete scale convention}
Fix $r\ge1$ and use $\Q^r$ with lexicographic order, the first coordinate being
the most significant. Embed it into the additive surreal group by
\begin{equation}\label{eq:embedding}
 (q_1,\ldots,q_r)\longmapsto
 q_1\omega^{r-1}+q_2\omega^{r-2}+\cdots+q_r.
\end{equation}
The powers of $\omega$ on the right are finite Conway powers. Uniqueness of
finite normal forms makes this an injective ordered additive map. Write
$\beta_j=\omega^{r-j}$ and $T_j=t^{\beta_j}$. In particular,
\[
 0<T_1<T_2^n<1\quad(n\in\N_{>0},\ r\ge2)
\]
when $T_1,T_2$ are successive scales in the case $r=2$; more generally every
earlier generator is smaller than every positive integer power of a later
one. For $r=2$, $T_1=\omega^{-\omega}$ and $T_2=\omega^{-1}$.

The order convention is not inferred from names such as \texttt{eps1} and
\texttt{eps2}; it belongs to the domain descriptor. Different conventions
would give different ordered embeddings, even if the same polynomial syntax
were used.

\subsection{Rational functions form a decidable ordered field}
\begin{theorem}[Effective rational monomial core]\label{thm:rationalcore}
Let
\[
 E_r=\RA(T_1,\ldots,T_r)\subset\No,
\]
with the embedding just specified. Finite rational-function descriptions give
an effective ordered field: arithmetic, equality, and sign terminate. The
same is true after admitting rational powers of the $T_j$, provided each
input uses finitely many such powers. The complexification supports decidable
field equality and conjugation.
\end{theorem}
\begin{proof}
Distinct integer exponent vectors give distinct Hahn monomials. A nonzero
polynomial in the $T_j$ with real algebraic coefficients therefore has a
nonempty finite normal form and cannot vanish. Thus the generators are
algebraically independent over $\RA$.

Polynomial arithmetic is effective because algebraic coefficient arithmetic
and equality are effective. A polynomial's sign is the sign of its coefficient
at the lexicographically least exponent in its finite support. For a fraction
$P/Q$, $Q\ne0$, equality is reduced to polynomial equality by cross
multiplication, and its sign is the product of the signs of $P$ and $Q$.
Addition, multiplication, and inversion of a nonzero fraction are finite
polynomial operations.

For rational exponents, choose a common positive denominator in each
coordinate for all exponents in the finitely many operands. Replace $T_j$
by $U_j^{N_j}$, where $U_j=t^{\beta_j/N_j}$. All the operands become rational
functions in the algebraically independent $U_j$. The same proof applies.
Finally represent complex elements by pairs of real elements, with the usual
formulas for arithmetic and conjugation.
\end{proof}
This core already admits $\omega$, $\omega^{-1}$,
$\omega^{-\omega}$, their rational powers, and rational combinations at
finitely many hierarchy levels. It does not contain arbitrary surreal
exponents such as every possible infinite normal form. Higher finite rank
can be added by an explicit order-preserving embedding of the old coordinates
into the new domain. In \eqref{eq:embedding}, preserving the old scales when
passing from rank $r$ to rank $r+1$ means prepending a zero coordinate.

\subsection{Adjoining all algebraic roots without expanding them}
Let $F_r$ be the relative real closure of $E_r$ inside $\No$: the elements of
$\No$ algebraic over $E_r$. It is a real closure because $\No$ is real closed.
Put $C_r=F_r[i]\subset\SC$.

\begin{theorem}[A practical algebraically closed target]\label{thm:algcore}
The field $F_r$ admits finite effective algebraic descriptions with decidable
order and equality, and $C_r$ admits finite effective descriptions with
decidable equality and field operations. Every finite-degree polynomial over
$C_r$ has all its roots in $C_r$. Finite polynomial systems over this domain
can be treated by effective algebraic methods, with the output representation
appropriate to their dimension.
\end{theorem}
\begin{proof}
Represent a real algebraic extension element by polynomial equations over
$E_r$ together with a root selector, for example a realizable Thom encoding.
A finite tuple of such descriptions can be handled in one real closure.
Real-closed-field quantifier elimination and algebraic sign determination
reduce questions about these descriptions to finitely many field operations
and sign tests in $E_r$. Those operations terminate by
\cref{thm:rationalcore}. This gives an effective presentation of the real
closure, including comparison and algebraic operations. Purely algebraic
quantifier-elimination methods based on Thom encodings are a published
foundation for this step \cite{PerrucciRoy}.

Pairs of real algebraic elements give an effective presentation of $C_r$;
complex equality is equality of both coordinates. Since $F_r$ is real closed,
$F_r[i]$ is algebraically closed. Consequently a polynomial over it factors
there, and no additional roots appear in a larger extension field. Algorithmic
root specifications can also be obtained by writing the polynomial equation
in real and imaginary coordinates and selecting its finitely many solutions.
For a polynomial system, the same real-algebraic encoding decides existential
solvability and supplies descriptions of its solution conditions. A
positive-dimensional solution set is represented by equations or a suitable
parametric/semialgebraic description, not by a finite list of points.
\end{proof}
The theorem is an algorithmic possibility statement for a carefully specified
field. It does not assert that Wolfram's native \texttt{Root} or
\texttt{Reduce} already accepts an arbitrary custom surreal head as its
coefficient domain. A practical implementation must supply the reductions,
ordered-field adapter, and root-description management.

\subsection{Why ordinary interval isolation is insufficient}
\label{subsec:nondense}
A subtle implementation trap is to assume that the ground ordered field is
dense in its real closure. It need not be. Take $E=\Q(t)$ with positive
infinitesimal $t$. Every nonzero rational function has integer valuation.
There is no $q\in E$ satisfying
\[
 \sqrt t<q<2\sqrt t.
\]
Indeed, any positive $q$ has leading form $ct^m$, $m\in\Z$, $c>0$. If
$m\le0$ it is larger than both endpoints; if $m\ge1$ it is smaller than both.
Consequently rational-function endpoints cannot separate the two roots
$\sqrt t$ and $2\sqrt t$ by a point between them.

Real-root indices and Thom encodings avoid this problem. Another approach
allows algebraic endpoints or enlarges the monomial field by fractional
powers. But naively porting an interval-bisection algorithm that assumes
rational density would not be a complete method. Similarly, a decimal
neighborhood at standard part zero cannot select one of several roots with
different infinitesimal displacements.

\subsection{What this core deliberately omits}
The countable union of these finite-rank algebraic cores is already useful,
but it is not the entire surreal field. In particular, algebraic closure does
not imply exponential closure, support closure under every countable sum,
or access to birthdays. Even an exactly named transcendental constant such as
$\pi$ is not part of the minimal algebraic coefficient domain. It can be added
through a larger coefficient backend, but the backend must declare what it
can decide about the new constants.

A symbolic parameter may instead remain algebraically independent over the
base. This permits rational-function algebra but is not an assertion that its
intended real value is independent of other named constants. The difference
between an indeterminate and a particular named number must survive
serialization, substitution, and simplification.

\section{Certified Hahn generators and finite dependency}
\label{sec:support}

\subsection{A support certificate is operational data}
The sources use strong summability: a family is summable when its union of
supports is well ordered and every exponent occurs in only finitely many
members. A software certificate should establish both properties. Merely
showing that every summand has positive valuation is insufficient: for
example, $\{1/n:n\ge1\}$ is a positive but non-well-ordered set of exponents.
Conversely, a well-ordered union alone does not justify infinitely many
contributions at a single exponent.

A useful support interface distinguishes the following requests. A
\emph{membership} operation determines whether an exponent lies in a support
bound. A \emph{coefficient} operation returns the actual coefficient, including
cancellation. A \emph{decomposition} operation lists the finitely many ways
an exponent can arise in a product or substitution. A \emph{leading-term}
operation additionally certifies that all smaller coefficients vanish. A
\emph{prefix} operation describes everything below a chosen cutoff. No one
of these operations should be inferred automatically from the existence of
the others.

For finite supports all of them are straightforward. For infinite supports,
certificates can be built compositionally: finite unions, translates,
Minkowski sums, positive finitely generated monoids, and finite products of
already certified constructors. An arbitrary well-order proof can justify
a mathematical expression without furnishing these operational services.

\subsection{An effective fixed-coefficient lemma for rational grids}
\begin{theorem}[Finite extraction from a positive rational grid]
\label{thm:grid}
Let $e_1,\ldots,e_m\in\Q^r$ be lexicographically positive, let $a\in\Q^r$,
and let $c:\N^m\to k$ be a total computable coefficient rule in an effective
field $k$. The monomial family
\begin{equation}\label{eq:gridseries}
 \left(c(\boldsymbol n)t^{a+\sum_jn_je_j}\right)_{\boldsymbol n\in\N^m}
\end{equation}
is strongly summable. At any specified $\gamma\in\Q^r$, its coefficient is
computable by a finite search, even if the map
$\boldsymbol n\mapsto a+\sum n_je_j$ is not injective.
\end{theorem}
\begin{proof}
The positive-monoid form of the Hahn--Neumann lemma makes the grid well
ordered and ensures finite fibers. We also give an effective bound.
There is a rational linear functional $\ell:\Q^r\to\Q$ such that
$\ell(e_j)>0$ for every $j$. One can find it by trying
\[
 \ell_M(x_1,\ldots,x_r)=M^{r-1}x_1+\cdots+Mx_{r-1}+x_r
\]
for successively larger positive integers $M$. For each lexicographically
positive $e_j$, its first nonzero coordinate eventually dominates the later
coordinates in this expression. Because there are finitely many $e_j$,
this search terminates.

Put $d=\gamma-a$. Every nonnegative integer solution of
$\sum n_je_j=d$ satisfies
\[
 \sum_jn_j\ell(e_j)=\ell(d).
\]
If $\ell(d)<0$, there are no solutions. Otherwise
\[
 0\le n_j\le\left\lfloor\frac{\ell(d)}{\ell(e_j)}\right\rfloor.
\]
Enumerate this finite box and test the vector equation exactly. The desired
coefficient is the finite sum of $c(\boldsymbol n)$ over its solutions.
The same bound proves directly the finite-contribution requirement for
strong summability. Zero coefficients can only decrease the support.
\end{proof}
The functional is an auxiliary device on the finite generator set. It is
not an order-preserving embedding of the whole lexicographic group into
$\Q$; no such embedding exists for $r>1$. Its use for fixed coefficients
therefore does not justify using it as a substitute valuation for arbitrary
cutoff comparisons.

This lemma supplies a concrete backend for a \texttt{GridSeries} object.
It also explains how to handle coincident exponent sums: sum their
coefficients rather than retaining multiple conflicting entries. A total
coefficient rule still needs a supplied termination argument, or membership
in a language where termination is built into the syntax.

\subsection{Known nonzero does not always make the leading term computable}
\begin{proposition}[A high-rank leading-term obstruction]
\label{prop:leading}
For total program-described two-scale grid series, even a promise that the
value is nonzero does not give a universal algorithm for its exact valuation
or its sign.
\end{proposition}
\begin{proof}
Use exponents in $\Q^2_{\rm lex}$, and let $a_n$ be the bounded-simulation
coefficients from \cref{thm:undecidable}, with halting steps indexed from $1$.
Define
\[
 B_M=t^{(1,0)}-\sum_{n\ge1}a_nt^{(0,n)}.
\]
This is always nonzero and actually has one or two terms. If $M$ never
halts, it is positive and has valuation $(1,0)$. If $M$ halts at step $n$,
it is negative and has valuation $(0,n)$. Either a sign algorithm or an
exact-valuation algorithm would decide halting.
\end{proof}
There is a useful contrast with a one-variable power series with a known
integer lower exponent bound. Under a nonzero promise, sequentially computing
its coefficients eventually finds the first nonzero coefficient. In higher
rank, infinitely many potential smaller exponents can precede a known nonzero
coefficient. An implementation must not generalize the rank-one search
argument without checking this issue.

\subsection{Inverses and positive recursions}
For a known nonzero series write
\[
 A=c\,t^a(1+E),\qquad c\ne0,\quad \supp(E)\subset\Gamma_{>0}.
\]
Then
\begin{equation}\label{eq:inverse}
 A^{-1}=c^{-1}t^{-a}\sum_{n\ge0}(-E)^n
\end{equation}
is an exact Hahn identity. For finite $E$, or a certified positive grid
representation of $E$, the support and finite decompositions can be handled
explicitly. If a general input has no computable leading normalization, the
formula is a mathematical identity rather than a ready algorithm. A package
can retain an exact rational-expression inverse without expanding it when
that representation is available.

Likewise, the supplied positive-operator formula
\[
 (I+T)^{-1}X=\sum_{n\ge0}(-T)^nX
\]
and its nonlinear preparation recursions have finite coefficient dependencies
controlled by support words \cite{SourceResearch}. To make them executable,
the implementation needs effective word decomposition and executable ordinary
coefficient operations. For a fixed-domain analytic division problem those
operations can themselves involve ordinary integrals or implicitly specified
functions. Such a coefficient is still representable by a well-defined
symbolic operator when no elementary closed form is available; it is not
thereby in the smallest decidable scalar field.

\section{Precision is a mathematical type, not a display setting}
\label{sec:precision}

\subsection{Finite truncations are not always possible}
Write
\[
 A=\widehat A+\Oh{\beta}
 \quad\text{to mean}\quad v(A-\widehat A)\ge\beta.
\]
This is a statement about a set of possible errors, not an equality in the
field after deleting a formal $O$-symbol. A stricter bound $v(\cdot)>\beta$
should be represented separately. The convention used in this article is
always the weak inequality unless a strict one is displayed.

\begin{proposition}[Finite-prefix criterion for rank-one grids]
\label{prop:finiteprefix}
For $a,e_1,\ldots,e_m,B\in\Q$ with every $e_j>0$, the grid
$a+\sum\N e_j$ contains only finitely many exponents below $B$, and those
exponents can be enumerated effectively. The corresponding statement is
false for general lexicographically ordered groups, even with one positive
generator.
\end{proposition}
\begin{proof}
If $a+\sum n_je_j<B$, then every $n_j<(B-a)/e_j$. A finite box therefore
contains all possibilities. For the counterexample use generator $(0,1)$
and cutoff $(1,0)$ in $\Q^2_{\rm lex}$. Every $(0,n)$ lies below that cutoff.
\end{proof}
Thus, in a workspace containing $1$ and $\omega$,
\[
 A=\frac1{1-t}=\sum_{n\ge0}t^n
\]
has no finite monomial truncation with error valuation at least $\omega$.
Indeed any finite polynomial misses some coefficient at a finite exponent,
and the first such difference has valuation below $\omega$. A compressed
answer such as the exact rational function remains possible. The obstruction
is to a \emph{finite term list}, not to every finite expression.

Even rank-one nondiscrete Hahn supports can fail the finite-prefix property.
The supplied example
\begin{equation}\label{eq:accumulating}
 A=\sum_{k\ge0}t^{1-1/(k+2)}
\end{equation}
has well-ordered support of order type $\omega$, with infinitely many terms
below $1$ \cite{SourcePolynomial,SourceResearch}. Hence ``rank one'' alone
is not enough; discreteness or a suitable positive-grid bound is needed.

\subsection{Four useful precision requests}
The API should distinguish a fixed coefficient, a finite number of leading
terms, a complete valuation prefix, and a finite multivariate jet in grid
coordinates. These requests are not interchangeable. A finite multivariate
jet is a quotient by a chosen monomial ideal in auxiliary variables; after
substituting the grid generators, its omitted monomials must be converted
into an explicit valuation error bound. That conversion may give a coarse
bound and may not describe an entire valuation prefix.

For higher-rank data, nested descriptions can be useful. A coefficient at one
outer scale may itself be an infinite series in a smaller-rank field. This
can compactly expose a term occurring after an order-type-$\omega$ block.
However, the nesting order must agree with the intended ordered group, and
operations must respect the resulting support conditions. A flat iterator
that only returns the next natural-number-indexed term cannot reach an
$\omega$th support position by a finite number of calls.

A precision request whose result cannot be a finite term list should return
a compressed certified subseries, an explicit unsupported-request result,
or an alternative finite jet chosen by the caller. It should never silently
stop after an arbitrary number of terms and claim the requested bound.

\subsection{Arithmetic with valuation error bounds}
\begin{proposition}[Basic certified precision rules]\label{prop:precision}
Suppose $A=\widehat A+e_A$, $B=\widehat B+e_B$, with
$v(e_A)\ge\alpha$ and $v(e_B)\ge\beta$. Then
\begin{align}
 v((A+B)-(\widehat A+\widehat B))&\ge\min\{\alpha,\beta\},\label{eq:addprecision}\\
 v(AB-\widehat A\widehat B)&\ge
 \min\{\alpha+v(\widehat B),\beta+v(\widehat A),\alpha+\beta\}.
 \label{eq:mulprecision}
\end{align}
If $\widehat A\ne0$, $\mu=v(\widehat A)$, and $\alpha>\mu$, then
$A\ne0$, $v(A)=\mu$, and
\begin{equation}\label{eq:invprecision}
 v(A^{-1}-\widehat A^{-1})\ge\alpha-2\mu.
\end{equation}
\end{proposition}
\begin{proof}
The sum error is $e_A+e_B$. The product error is
$e_A\widehat B+\widehat A e_B+e_Ae_B$. Apply the valuation inequalities.
Under the last hypothesis, the leading term of $\widehat A$ cannot be canceled
by $e_A$, so $v(A)=\mu$. Finally
\[
 A^{-1}-\widehat A^{-1}=-\frac{e_A}{A\widehat A},
\]
whose valuation is at least $\alpha-2\mu$.
\end{proof}
These formulas expose genuine precision loss through inversion of a small
quantity. For instance, an error of valuation $10$ in a quantity of valuation
$3$ gives an inverse error bound of only $4$. It would be incorrect to
preserve the same absolute cutoff through division.

A zero test on a truncated series may return ``zero to the supplied
precision,'' but that is not exact zero. If a nonzero leading coefficient is
known below the error cutoff, sign and leading valuation are certified. If
all retained terms cancel, more information is required. Exact identity
tracking can still prove $A-A=0$ even when only an approximate expansion of
$A$ has been computed; treating the two appearances as unrelated error balls
would lose this information unnecessarily.

\subsection{Numerical coefficient precision is a second axis}
Surreal scale precision and ordinary numerical precision are independent.
The coefficient of $t^{1/2}$ may be known to 100 binary digits while the
series is known only through exponent $3$. Increasing coefficient precision
does not discover an omitted $t^4$ term. Conversely, producing more exact
rational coefficients does not yield an ordinary floating-point
approximation to an infinite surreal value.

For nonalgebraic ordinary coefficients, validated real or complex enclosures
can sometimes certify nonvanishing and sign. If an enclosure contains zero,
one must refine it or retain an unresolved status. Heuristically small is not
exactly zero. An ordinary approximate coefficient also cannot be silently
recast as its nearby rational value in an exact-number constructor.

\subsection{Why these are not fine-topological limits}
The full-class fine topology in the supplied analysis manuscript uses every
positive surreal tolerance. Every set-indexed convergent net in that topology
is eventually equal to its limit \cite{SourceAnalysis}. A sequence of finite
truncations should therefore not be advertised as converging in that topology.
Within a chosen rank-one valued field, a sequence of truncations can converge
in its valuation topology; the same words refer to a different structure.

In particular, writing \texttt{N[surreal,p]} requires a declared meaning.
Reasonable choices are a certified coefficient approximation, a valuation jet,
or a scale-aware structured result. Converting all infinities to native
\texttt{Infinity} and all infinitesimals to zero is neither an injective
representation nor a field homomorphism.

\section{Polynomial algebra and root-cluster algorithms}
\label{sec:polynomials}

\subsection{Finite algebra over a declared coefficient domain}
Over an effective field, polynomial degree is an ordinary finite integer.
Long division terminates by degree descent. The Euclidean algorithm gives
gcds and B\'ezout identities; squarefree decomposition uses the derivative;
resultants and discriminants are finite determinants; and matrices admit
exact row reduction whenever pivot zero tests are available. These algorithms
apply to the core of \cref{sec:core} without requiring an expansion of every
coefficient into an infinite Hahn normal form.

For multivariate polynomials, Buchberger-type ideal algorithms depend on
finite monomial orders and effective coefficient operations. A completed
Gr\"obner basis can certify ideal membership or inconsistency by a finite
identity. A zero-dimensional quotient has a finite basis and multiplication
matrices; a positive-dimensional solution set needs a different output shape.
The mathematical existence of a root over an algebraically closed Hahn field
is not a promise to list a positive-dimensional family.

If coefficient equality is only partial, the polynomial routines must become
partial or conditional as well. An unresolved leading coefficient can change
the degree. An unresolved pivot can change matrix rank. Testing a coefficient
with an ordinary numerical heuristic and proceeding as though the result were
an exact algebraic fact is unsound.

\subsection{The initial polynomial is a valuable finite output}
The supplied polynomial manuscript gives a particularly useful computational
interface. Write
\[
 P(a+Y)=\sum_{j=0}^n b_jY^j,\qquad
 \lambda=\min_{b_j\ne0}\{v(b_j)+j\rho\}.
\]
Define the nonzero ordinary residual polynomial
\begin{equation}\label{eq:initial}
 I_{a,\rho}(P)(X)=\st\!\left(t^{-\lambda}P(a+t^\rho X)\right).
\end{equation}
Its computation uses finite translation, exact coefficient valuations and
leading coefficients, exponent comparison, and residue arithmetic. It need
not find any full root expansion.

For a complete root list $\alpha_i$ in a containing algebraically closed
workspace, the source proves
\begin{equation}\label{eq:initialfactor}
 I_{a,\rho}(P)(X)=C
 \prod_{v(\alpha_i-a)\ge\rho}
 \left(X-\st\frac{\alpha_i-a}{t^\rho}\right),\qquad C\ne0.
\end{equation}
Consequently its degree counts the roots in the closed valuation ball,
its order at zero counts the roots in the open ball, and the multiplicity
of a residual root $c$ counts the roots in the corresponding residue-direction
ball. These are exact finite conclusions even when the represented roots
have infinite supports \cite{SourcePolynomial}.

The required valuations are effective for rational monomial coefficients by
\cref{thm:rationalcore}. For algebraic or more general coefficients, an
additional valuation/leading-coefficient service is needed. Decidable field
equality alone should not be confused with an already implemented expansion
engine. Algebraic root descriptors can remain valid while that service is
still being developed.

\subsection{Newton profiles do not require a real-valued plot}
At center zero the profile is
\[
 \phi_P(\rho)=\min_j\{v(a_j)+j\rho\}.
\]
Candidate breakpoints are the finitely many values
\[
 \rho=\frac{v(a_j)-v(a_k)}{k-j}\quad(j\ne k).
\]
At a breakpoint, the difference between the largest and smallest active
indices is the number of nonzero roots at that valuation. Exact division by
ordinary integers in a divisible exponent group is enough; the exponents
need not be embedded in $\R$ \cite{SourcePolynomial}.

A straightforward implementation enumerates candidates, sorts them by the
ordered-group comparator, and evaluates all affine expressions at each
candidate. A lower-hull implementation can reduce repeated work. In either
case, drawing an ordinary Euclidean Newton polygon is only a visualization
for suitable rank-one inputs. For exponents such as $1$ and $\omega$, a
finite algebraic profile is more faithful than a misleading numerical plot.

\begin{example}[The cubic from the supplied polynomial manuscript]
\label{ex:cubic}
Let $P(z)=z^3-t^2z-t^5$. Its coefficient profile is
\[
 \min\{5,2+\rho,3\rho\}.
\]
The active indices at $\rho=1$ are $1,3$, and those at $\rho=3$ are $0,1$.
Thus two roots have valuation $1$ and one has valuation $3$. The residual
polynomials are
\[
 I_{0,1}(P)=X^3-X,\qquad I_{0,3}(P)=-X-1.
\]
The nonzero residual roots select leading terms $t,-t,-t^3$. After these
simple residual branches have been selected, ordinary coefficient recursion
gives
\begin{align*}
 \alpha_+&=t+\tfrac12t^3-\tfrac38t^5+\Oh7,\\
 \alpha_-&=-t+\tfrac12t^3+\tfrac38t^5+\Oh7,\\
 \alpha_0&=-t^3-t^7+\Oh{11}.
\end{align*}
These finite displays do not replace the exact root descriptors. Rather,
they are certified expansions attached to those descriptors.
\end{example}

\subsection{Hensel lifting, multiple clusters, and algebraic fallback}
When a monic polynomial has pairwise coprime ordinary residual factors,
the source constructs unique lifted cluster factors. At each positive
exponent the correction is obtained by the same finite linear inverse, and
the support lies in the positive monoid generated by the input correction
support \cite{SourcePolynomial}. With effective support decomposition this
is an implementable coefficient algorithm.

A multiple residual root is different. For $z^m-t$, every root reduces to
zero, and the first useful scale is $t^{1/m}$. A routine which only attempts
ordinary simple-root Newton iteration at zero will fail. It should split or
rescale the cluster, enlarge the exponent group when required, or return an
exact algebraic object with an unresolved expansion. The algebraic fallback
is particularly valuable: it separates exact solvability from the cost of
finding a convenient normal form.

There is no unrestricted guarantee that a natural-number sequence of
successive corrections becomes cofinal in every desired higher-rank cutoff.
The source's general lifting proof is a well-founded support recursion. For
rank-one algebraic series, classical Newton--Puiseux methods provide a much
more specialized effective expansion mechanism; that specialization should
be identified rather than extrapolated to arbitrary program-described Hahn
coefficients \cite{Poonen,SourcePolynomial}.

\subsection{How much precision is enough for a root assertion?}
For monic degree-$n$ polynomials with finite coefficients, the source proves
that coefficient errors of valuation at least $\varepsilon>0$ permit a
matching of root multisets with
\begin{equation}\label{eq:rootmatch}
 v(\alpha_i-\beta_i)\ge\varepsilon/n.
\end{equation}
The exponent loss is sharp: compare $z^n$ with $z^n-t^\varepsilon$.
A software request for roots to valuation $\rho$ may therefore require
coefficient information through approximately $n\rho$ near an unresolved
multiple cluster. Ignoring multiplicity would overstate the precision.

For a squarefree integral polynomial, put
\[
 d_i=v(P'(\alpha_i)),\qquad
 \delta_i=\max_{j\ne i}v(\alpha_i-\alpha_j).
\]
The source gives the sharper sufficient condition
$\varepsilon>\max_i(d_i+\delta_i)$, with
\[
 v(\beta_i-\alpha_i)\ge\varepsilon-d_i.
\]
An easily stated coefficient-level sufficient condition is
$\varepsilon>v(\Disc P)$ \cite{SourcePolynomial}. These are suitable
certificates for stable cluster labels, not merely heuristic tolerances.
The strictness matters:
\[
 z(z-t^a)+\frac{t^{2a}}4=(z-t^a/2)^2
\]
shows a collision exactly at the threshold. Before applying these estimates
to infinite root locations, the polynomial must be normalized to the stated
finite-coefficient setting and the error transformed with it.

\subsection{Other implementable finite consequences}
Once the initial polynomial and the requisite residue factorization are
available, the source's theorems also give exact counts in polynomial image
balls and critical-point counts in occupied root balls. A ball containing
$k\ge1$ root multiplicities contains $k-1$ derivative-root multiplicities.
These routines should retain the occupation hypothesis: a ball containing no
roots of $P$ can still contain a root of $P'$. They should also distinguish
root counts from explicit root locations.

Order-geometric predicates, such as a polynomial inequality on a circle with
surreal radius, can be encoded in real-closed-field arithmetic at fixed finite
degree. This gives a route through the core's sign-determination machinery.
It is distinct from valuation-ball analysis, and it does not create a
fine-continuous contour integration routine.

\section{Exact functions: formal germs, ordinary lifts, and coherent data}
\label{sec:functions}

\subsection{Formal evaluation at an infinitesimal}
Let $k\subseteq\C$ be an admitted coefficient field and let
$A(X)=\sum_{n\ge0}a_nX^n$ be a specified formal power series over $k$.
For an infinitesimal $\epsilon$ with a certified well-ordered positive
support, the evaluation
\begin{equation}\label{eq:formaleval}
 A(\epsilon)=\sum_{n\ge0}a_n\epsilon^n
\end{equation}
is strongly summable. Indeed, the support is contained in
$(\supp\epsilon)^*$, and finite word multiplicity gives finitely many
contributions at each exponent. No ordinary radius of convergence is required
\cite{SourceAnalysis}.

Thus both $\sum t^n/n!$ and $\sum n!t^n$ have exact Hahn interpretations.
Only the first is the Taylor series of an ordinary holomorphic function near
zero. The second may be stored as a formal-series value, with a recurrence
and support certificate. Calling both objects ordinary analytic germs would
erase a distinction on which the supplied coherence theory depends.

For arbitrary surreal coefficients $a_n$, positivity of $v(\epsilon)$ is
not enough: the coefficient valuations may decrease too rapidly. The source
proves existence of a sufficiently small scale for any set-sized coefficient
family, possibly outside a previously fixed workspace. For software, a usable
scale or an effective bound on the coefficient supports must be supplied or
computed. The mathematical existence of such a scale is not a universal
algorithm for choosing it from an arbitrary program.

\subsection{Ordinary holomorphic functions lift canonically on their domains}
For an ordinary holomorphic function $f$ on $U\subseteq\C$, define
\begin{equation}\label{eq:lift}
 f\sh(c+\epsilon)=\sum_{n\ge0}\frac{f^{(n)}(c)}{n!}\epsilon^n,
 \qquad c\in U,\quad \epsilon\in\mm.
\end{equation}
This is the scalar lift in the supplied manuscripts. It preserves ordinary
holomorphic identities and admissible compositions. The domain is
$U\sh=\{c+\epsilon:c\in U,\epsilon\in\mm\}$, not every surcomplex number
whose printed formula resembles an ordinary argument.

Exact symbolic values of $\exp$, $\sin$, $\cos$, inverse trigonometric
branches, logarithm branches, and ordinary special functions can be admitted
through this mechanism wherever the relevant ordinary function is
holomorphic and specified. At an ordinary pole, a meromorphic germ can
instead be factored into a finite negative power and a holomorphic germ.
Evaluation then requires a nonzero displacement. At a branch point,
ramification or a chosen algebraic parameter is additional data.

For example, the expression
\[
 \sin\sh(\pi+t)=-\sin\sh(t)
\]
is exact once the ordinary constants and functions are admitted. But a
minimal $\RA$-coefficient implementation does not thereby gain a complete
decision procedure for every ordinary constant generated by arbitrary special
function evaluation. A value constructor may be more general than its
simplifier.

\subsection{The Hahn-coherent layer follows the sources literally}
A Hahn-coherent datum on a fixed ordinary domain is
\[
 F=\sum_{\gamma\in S}f_\gamma t^\gamma,
 \qquad f_\gamma\in\OO(U),\quad S\text{ well ordered}.
\]
Its value at $c+\epsilon$ is
\begin{equation}\label{eq:coherenteval}
 F\sh(c+\epsilon)=
 \sum_{\gamma\in S}\sum_{n\ge0}
 t^\gamma\frac{f_\gamma^{(n)}(c)}{n!}\epsilon^n.
\end{equation}
The source support bound is $S+(\supp\epsilon)^*$. Coefficient functions
must share the stated ordinary domain. Pointwise knowledge that selected
values are infinitesimal does not substitute for this uniform support and
domain condition \cite{SourceAnalysis,SourceResearch}.

A software representation can therefore contain a domain object, a support
certificate, and a coefficient-function provider. Finite support with rational
or explicitly holomorphic coefficient functions is a particularly accessible
initial implementation. Infinite coherent data need further operational
certificates, just as infinite scalar series do.

The source's Hartogs and gluing theorems are mathematically useful, but they
should not be exposed as unconditional algorithms that discover a global
support or decide arbitrary holomorphic identities. To invoke them
computationally, the input must contain the relevant function descriptions
and hypothesis certificates, or a restricted syntax in which those
hypotheses can be checked.

\subsection{Composition has a domain and support contract}
For $F=f_0+E$ with nonnegative support, $E$ of positive support, and
$f_0(U)\subseteq V$, the source defines
\begin{equation}\label{eq:composition}
 H\circ F=
 \sum_{\delta,n}\frac{h_\delta^{(n)}\circ f_0}{n!}\,t^\delta E^n,
 \qquad H=\sum_\delta h_\delta t^\delta\in\HH(V).
\end{equation}
The support is contained in $\supp H+(\supp E)^*$.
An executable implementation needs a verified ordinary-domain inclusion,
ordinary derivatives and compositions, and effective support bookkeeping.
A symbolic domain obligation can remain in a conditional result when it
cannot be resolved automatically.

Blind substitution fails. The geometric expansion
$\sum_{n\ge0}t^nz^n$ is coherent on the finite plane and equals
$(1-tz)^{-1}$ there. Substituting $z=t^{-2}$ into that expansion produces
exponents $-n$, which are not well ordered. The rational expression itself
still has a value at $t^{-2}$, but its valid expansion is now
\[
 \frac1{1-t^{-1}}=-\frac{t}{1-t}=-\sum_{n\ge1}t^n.
\]
At $z=t^{-1}$ there is a pole. Thus a failed series substitution can require
re-expansion, a different representation, or rejection, depending on the
actual argument.

\subsection{Implicit functions and local root specifications}
For a formal system $\Phi(X,Y)=0$ with specified initial point and invertible
Jacobian in $Y$, the coefficient of each new homogeneous degree is determined
by one fixed finite matrix inverse. When the coefficient operations and the
input generators are effective, the solution's finite jets are computable.
This is a natural exact object:
\texttt{ImplicitSeries[equations, variables, initialData, certificate]}.
The equation and uniqueness certificate denote the whole solution; the
computed jet is only a view of it.

In the coherent category, the sources' preparation and matrix-division
constructions preserve one ordinary domain and explicit positive-monoid
support. They can therefore serve as exact lazy constructors. Their
implementation is more demanding than the formal degree recursion because
the coefficients are ordinary functions on that domain, and division must
preserve holomorphy there. A general-purpose CAS need not find closed forms
for all such coefficient functions in order to retain them exactly as
operator expressions.

\subsection{A concrete nonpolynomial algebraic cluster}
The supplied example $\sin^2z-t$ has two zeros in the zero monad:
\[
 z_\pm=\pm\arcsin\sh(\sqrt t)
 =\pm\left(t^{1/2}+\frac{t^{3/2}}6+
 \frac{3t^{5/2}}{40}+\frac{5t^{7/2}}{112}+\cdots\right).
\]
The exact prepared cluster polynomial is
\[
 P(z)=z^2-\bigl(\arcsin\sh(\sqrt t)\bigr)^2.
\]
It is polynomial in $z$, but its coefficients lie in an analytic extension
of the minimal rational monomial field. A package should not label the roots
``algebraic over the original coefficient field'' merely because they occur
in a prepared polynomial over a newly enlarged field. The ambient coefficient
domain is part of every algebraicity claim.

\section{Exponentials, logarithms, and transcendental scale semantics}
\label{sec:transcendental}

\subsection{Canonical real exponentiation is a valid exact constructor}
The Gonshor exponential $\exp_{\No}$ and its positive-argument inverse
$\log_{\No}$ are established structures on the surreals
\cite{Gonshor,vdDE}. A finite expression such as
\[
 \exp_{\No}(\omega),\qquad
 \log_{\No}(\omega+1),\qquad
 \exp_{\No}(\omega+\omega^{-1})
\]
therefore has an exact meaning. The fact that a chosen implementation cannot
yet expand one of these expressions does not make its denotation inexact.
For the last example the exact group law gives
\[
 \exp_{\No}(\omega+\omega^{-1})
 =\exp_{\No}(\omega)\sum_{n\ge0}\frac{\omega^{-n}}{n!}.
\]
This is useful as a structured monomial-scale expression even without
normalizing the first factor in the smallest Hahn backend.

Supporting the constructor is easier than supporting a complete zero,
comparison, or normal-form algorithm for every expression built from it.
A practical extension should have a documented rewrite theory, a declared
transseries grammar when used, and an unresolved outcome for identities
outside that theory. It should also make positive-real powers
$\exp_{\No}(b\log_{\No}a)$ distinct from the Conway monomial notation
of \cref{subsec:omega}.

\subsection{Infinite imaginary arguments require an extension policy}
\label{subsec:phase}
The supplied analysis manuscript proves a concrete nonuniqueness result.
For a real surreal $y$, let $\operatorname{fin}(y)$ retain its nonnegative
$t$-exponent terms, and let $\ell(y)$ be its coefficient at $t^{-1}=\omega$.
For an ordinary real $\alpha$, put
\[
 \theta_\alpha(y)=\operatorname{fin}(y)+\alpha\ell(y),
\]
and define
\begin{equation}\label{eq:phaseexp}
 E_\alpha(x+iy)=\exp_{\No}(x)
 \left(\cos\sh\theta_\alpha(y)+i\sin\sh\theta_\alpha(y)\right).
\end{equation}
According to the source, these are global fine-analytic group homomorphisms
extending the ordinary complex exponential and satisfying $E_\alpha'=E_\alpha$.
Nevertheless,
\[
 E_0(i\omega)=1,\qquad E_\pi(i\omega)=-1.
\]
Thus agreement on ordinary complex numbers, the group law, and fine-local
analyticity do not select a unique value at $i\omega$
\cite{SourceAnalysis}.

A CAS should consequently not assign an unqualified value to
\texttt{SurcomplexExp[I omega]} merely by extending ordinary Taylor syntax.
It may require an explicit extension identifier, implement a documented
stronger theory, or leave the operation outside its admitted domain.
This statement does not assert that no other published convention can be
chosen. It says that the source's stated local axioms do not make that choice.

Even formula \eqref{eq:phaseexp} is not automatically an effective algorithm
on every exact Hahn generator: extracting the whole finite part can be a
nontrivial support-cut operation. On finite sparse data and suitably
certified structured data, it is straightforward. On an unrestricted program
representation, its operational availability must be declared.

\subsection{Logarithms need branch and category metadata}
For a nonzero surcomplex $z$, the source's global fine-analytic logarithm is
specified by
\[
 r=|z|,\quad u=z/r,\quad u_0=\st(u),\qquad
 \mathcal L(z)=\log_{\No}r+i\Arg(u_0)+\Log_{\rm inf}(u/u_0).
\]
The ordinary argument convention is part of the definition. This logarithm
is a right inverse of the source's $E_\alpha$, but is not a coherent
primitive of $1/z$ on the entire thickened ordinary punctured plane
\cite{SourceAnalysis}. A local ordinary logarithm branch, a global fine-local
logarithm, and the real positive logarithm are therefore different function
types even when their printed name is \texttt{Log}.

A safe simplifier may use $\exp_{\No}(\log_{\No}x)=x$ for positive real
surreal $x$. It may use a chosen surcomplex right-inverse identity when the
chosen constructions justify it. It must not use an unrestricted
$\log(\exp z)=z$, or cancel arbitrary complex powers without branch
hypotheses. Returning a set of roots is often preferable to imposing a
hidden principal-root policy on a nonstandard complex plane.

\subsection{Transseries are a natural but explicitly selected extension}
Grid-based exponential--logarithmic transseries provide ordered scale
hierarchies and a substantial algebraic and differential apparatus
\cite{vdH}. Berarducci and Mantova study distinguished surreal-function
classes with composition and Taylor behavior, and separately construct a
compatible surreal derivation \cite{BMFunctions,BMDerivations}. These results
are relevant precedents for a transserial backend. They do not justify
substituting arbitrary surreals into every transseries expression without
checking the chosen framework's hypotheses.

An implementation might admit a finite expression grammar of real algebraic
constants, an infinite variable, ordered exponential and logarithmic scales,
and certified positive-grid expansions. Operations would either return an
object in the same grammar, enlarge the finite scale list with a checked
embedding, or return a symbolic expression whose more ambitious operations
are not yet available. The exact closure rules should be documented for that
grammar. The entire mathematical field of grid-based transseries is not a
finite-input format when arbitrary coefficient arrays are allowed.

Oscillatory functions deserve a separate layer. An expression such as
$\sin x$ at large positive real $x$ does not have an eventual sign in the
same sense as a nonzero ordered logarithmic--exponential transseries.
Treating oscillation as just one more positive ordered monomial would lose
its essential behavior. For infinite surreal arguments, the extension
convention from the preceding subsections must also be addressed.

\subsection{Asymptotic data do not determine exact values by themselves}
The real functions $f(x)$ and $f(x)+e^{-x}$ have the same asymptotic expansion
in inverse powers of $x$ whenever the first expansion exists: $e^{-x}$ is
smaller than every such power as $x\to+\infty$. Under a selected compatible
surreal extension, however, the additional term at $x=\omega$ is a nonzero
infinitesimal. Therefore knowing all algebraic orders of an asymptotic
expansion does not in itself select one exact value.

Similarly, the strong Hahn value $\sum n!t^n$ is not automatically a Borel
sum, a lateral resummation, or the value of a particular analytic function
at $1/t$. Such identifications require a summation and extension theorem,
and can carry sector and Stokes data. Costin and Ehrlich's work on surreal
extension and integration supplies a relevant published direction, but no
blanket identification with the coefficientwise contour theory of the
supplied manuscripts is assumed here \cite{CostinEhrlich,SourceAnalysis}.

\subsection{All-scale rigidity is an interface boundary}
The supplied analysis manuscript's all-scale Hahn-entire functions are
precisely polynomials, and its all-scale Hahn-meromorphic functions are
rational \cite{SourceAnalysis}. The quantifier ranges over all surreal scales
and uses the manuscript's uniform ordinary-coefficient chart condition.
A package that admits a global exponential should therefore not also label
it ``all-scale Hahn-entire'' in that sense. The correct conclusion is a
separation of categories, not a prohibition on transcendental constructors.

\section{Differentiation, integration, and three different variables}
\label{sec:calculus}

\subsection{A coordinate derivative fixes surreal coefficients}
For a polynomial or coherent function in a variable $z$,
\[
 F(z)=\sum_\gamma f_\gamma(z)t^\gamma,
 \qquad D_zF=\sum_\gamma f_\gamma'(z)t^\gamma.
\]
Here every $t^\gamma$ is a coefficient, so $D_z(t^\gamma)=0$. This is the
coordinate derivative used in the supplied analytic manuscripts. It permits
ordinary symbolic differentiation, Taylor jets, Jacobians, implicit
differentiation, and residues in $z$.

For a formal power series in an auxiliary indeterminate $T$, the formal
derivative $D_T$ instead satisfies $D_TT=1$. Evaluating that formal series at
a fixed surreal $t$ produces a number, not an ordinary variable on which
$D_z$ should suddenly act. A package should record whether a symbol is a
coordinate, a formal parameter, or a designated surreal constant.

\subsection{A surreal-field derivation is additional structure}
A surreal derivation, such as a chosen Berarducci--Mantova derivation
normalized by $\partial\omega=1$, is a derivation of the number field itself
\cite{BMDerivations}. Then
\[
 t=\omega^{-1}\quad\Longrightarrow\quad \partial t=-t^2,
\]
whereas $D_TT=1$ and $D_zt=0$. For rational $q$,
\[
 \partial(t^q)=-q\,t^{q+1}.
\]
This last formula is a rational-power field identity; it should not be
extended to every surreal exponent of the Conway omega-map by treating that
exponent as an ordinary constant and applying an unjustified elementary
power rule.

Accordingly, a proposed interface should have distinct operations such as
\texttt{CoordinateD}, \texttt{FormalD}, and
\texttt{SurrealDerivation[..., derivationID]}. Automatic use of native
\texttt{D} is reasonable only after the chosen coordinate semantics have
been established.

\subsection{Integration can be exact without an elementary antiderivative}
In a formal power-series ring of characteristic zero, termwise integration
with zero constant term is effective when the coefficients are effective.
A Laurent term $T^{-1}$ requires adjoining a logarithm; the integral need
not remain Laurent or Puiseux. In a selected differential field, elementary
integration algorithms or differential-equation objects may supply broader
representations, but closure of the whole surreal field under a chosen
antiderivation is not a universal finite algorithm on arbitrary inputs.

For Hahn-coherent functions the sources define a different operation:
\begin{equation}\label{eq:contour}
 \int_\gamma^{H}F(z)\,dz
 =\sum_\lambda t^\lambda\int_\gamma f_\lambda(z)\,dz,
\end{equation}
where $\gamma$ is an ordinary path and the integral is coefficientwise.
On a simply connected ordinary domain, normalized primitives are constructed
coefficientwise. On a multiply connected domain, periods matter
\cite{SourceAnalysis}.

This gives an exact operator representation even when an ordinary coefficient
integral has no elementary closed form. To turn it into a computable scalar
coefficient one needs an ordinary exact or validated integration method for
that coefficient class. Exchanging a general numerical quadrature result for
an exact coefficient would change the contract. The contour in
\eqref{eq:contour} is not a fine-continuous real-parameter path through the
surcomplex class.

\subsection{Some familiar operations do not transfer literally}
Ordinary \texttt{Floor}, \texttt{Ceiling}, and integer rounding usually target
the standard integers. There is no standard integer between $\omega-1$ and
$\omega$, and no ordinary integer floor of $\omega$. A chosen surreal integer
part, for example an omnific-integer convention, is a different operation and
needs its own definition. Similarly, ordinary decimal precision does not
supply a finite float for an infinite value.

Differential equations require the same care. A fine-local solution of
$f'=f$ with a value at zero need not be globally unique in the entire
fine-local category, as the supplied exponential family shows. A local
formal initial-value problem can still be uniquely solved under the usual
coefficient-recursion hypotheses. A solver must state the category in which
its uniqueness claim is made.

\section{Finite quotient algebras: compute invariants before expanding roots}
\label{sec:quotients}

\subsection{Remainders, traces, and residues avoid branch proliferation}
Let $P\in K[z]$ be monic of degree $n$, where $K$ is an effective surcomplex
coefficient domain, and set $A_P=K[z]/(P)$. Polynomial remainders represent
elements of this $n$-dimensional algebra. Multiplication matrices give traces,
norms, characteristic polynomials, and resultants without choosing individual
roots. The source's residue functional is
\begin{equation}\label{eq:residuefunctional}
 \lambda_P(H)=[z^{n-1}](H\bmod P),
\end{equation}
and its trace identity is
\[
 \Tr(m_H)=\lambda_P(P'H).
\]
Both are finite exact coefficient calculations \cite{SourcePolynomial}.

The pairing $(G,H)\mapsto\lambda_P(GH)$ has determinant
$(-1)^{n(n-1)/2}$ in the power basis, even at repeated roots. The trace pairing,
in contrast, has determinant $\Disc P$ and can degenerate. This is an
important algorithmic distinction: a residue calculation by remainders can
remain regular through a collision even when formulas using individual
weights $1/P'(\alpha_i)$ develop apparent singularities.

\subsection{An analytic observable in a polynomial quotient}
For $P=z^2-t$, the relation $z^2=t$ gives the formal reduction
\[
 e^z\bmod P=A(t)+zB(t),
\]
where
\[
 A(t)=\sum_{n\ge0}\frac{t^n}{(2n)!},\qquad
 B(t)=\sum_{n\ge0}\frac{t^n}{(2n+1)!}.
\]
These are strongly summable at infinitesimal $t$. Hence
\begin{equation}\label{eq:residueexample}
 \lambda_P(e^z)=B(t)
 =1+\frac t6+\frac{t^2}{120}+\frac{t^3}{5040}+\cdots.
\end{equation}
Equivalently, the source's normalized contour integral of
$e^z/(z^2-t)$ equals $\sinh\sh(\sqrt t)/\sqrt t$. The individual moving-pole
residues are infinite, but their sum is finite \cite{SourceAnalysis}.

The quotient computation avoids both a branch choice for the intermediate
square root and subtracting two infinite leading residue weights. This
suggests a useful CAS policy: compute symmetric invariants, remainders, and
finite-algebra traces before requesting full labeled root expansions.
For more general analytic observables, an exact remainder constructor may
be supported by the coherent division theorem and a certificate, rather than
by a finite polynomial reduction alone.

\subsection{Coupled systems and the ordinary-domain contract}
For the supplied system
\[
 F_1=x^2-ty,\qquad F_2=y^2-tx,
\]
the quotient has basis $(1,x,y,xy)$ and dimension four. Its roots are
$(0,0)$ and $(t\zeta,t\zeta^2)$ for $\zeta^3=1$. The residue pairing has
nonzero constant determinant despite infinite individual evaluation weights
\cite{SourceResearch}.

More generally, that manuscript treats coupled coherent systems
$F_i=w_i^{m_i}+E_i$ with positive-support errors on one ordinary product
domain. It constructs a finite-free algebra of rank $\prod_i m_i$ and
normal forms with support in $S+E^*$. This can motivate a lazy
\texttt{CoherentQuotientAlgebra} constructor whose multiplication matrices
are computed coefficientwise. It is not a generic theorem for arbitrary
analytic systems or arbitrary ordinary leading complete intersections; the
source explicitly delimits those extensions.

The source also distinguishes the proved multidimensional residue pairing
from a general identification of every localized functional with a chosen
Grothendieck-residue theory. A software implementation should retain that
boundary. Verifying the Jacobian trace identity in a worked example does
not install it as a universal rewrite rule for all coherent quotients.

\section{What Wolfram Language already supplies, and what an extension must add}
\label{sec:wolfram}

\subsection{A documentation-based assessment}
The following assessment uses official documentation inspected for this
article on September 21, 2026. It identifies reusable facilities, not a claim
that every function automatically accepts an arbitrary user-defined field.
No documented native object with all of the arbitrary-Hahn and surreal
contracts developed here was identified in the inspected material. This is
not an exhaustive survey of third-party packages.

\begin{longtable}{@{}P{0.24\textwidth}P{0.33\textwidth}P{0.35\textwidth}@{}}
\toprule
Facility & What it contributes & What still needs a surreal adapter\\
\midrule\endhead
\texttt{Root}, \texttt{AlgebraicNumber}, \texttt{RootReduce} & Exact ordinary
algebraic coefficients and implicit-root precedents. & Ordered monomial
domains, non-Archimedean selectors, and custom coefficient semantics.\\[0.5em]
\texttt{SeriesData}, \texttt{SeriesCoefficient} & Finite power/Puiseux series
views, coefficient extraction, and order tracking. & Arbitrary ordered
exponent groups, infinite-prefix handling, and Hahn support certificates.\\[0.5em]
\texttt{Series}, \texttt{InverseSeries}, composition & Ordinary local
expansions and reversion, including supported fractional or logarithmic
forms. & A proof that the requested surreal evaluation is admissible and
that the chosen branch is the intended one.\\[0.5em]
\texttt{Asymptotic}, \texttt{AsymptoticSolve} & Ordinary asymptotic
approximations and local equation expansions. & Exact Hahn denotation,
beyond-all-orders data, and selected surreal extension semantics.\\[0.5em]
\texttt{DifferentialRoot} & A finite ordinary-function representation using
linear differential equations and initial conditions. & Strong evaluation,
support and singular-point certificates, and coefficient-domain closure.\\[0.5em]
Polynomial and matrix routines & Reusable finite algebra and symbolic
coefficient manipulations. & Exact zero tests and specialization into the
declared ordered field; not a generic automatic custom-field interface.\\[0.5em]
\texttt{UpValues}, custom heads, \texttt{Inactive} & Extension mechanisms and
controlled symbolic syntax. & Mathematical validity, termination policy,
branch conventions, and deliberate numerical behavior.\\
\bottomrule
\end{longtable}
The representation details are documented in
\cite{WLRoot,WLAlgebraic,WLRootReduce,WLSeriesData,WLSeriesCoefficient,WLSeries,
WLInverseSeries,WLAsymptotic,WLAsymptoticSolve,WLDifferentialRoot,
WLUpValues,WLInactive}. The proposed adapters in the last column are not
native-feature claims.

\subsection{Do not understate the current scope of \texttt{Root}}
Current documentation distinguishes indexed polynomial-root representations
from neighborhood-based representations, and includes exact roots of certain
transcendental equations and systems \cite{WLRoot}. Thus the assertion
``\texttt{Root} can only represent algebraic numbers over $\Q$'' would be
incorrect. Its ordinary algebraic subset remains particularly useful for an
exact coefficient backend.

None of those documented ordinary numerical neighborhood mechanisms, by
itself, chooses between roots whose standard parts are identical but whose
surreal displacements differ. Nor does an ordinary symbolic parameter in
\texttt{Root} automatically acquire the ordering of an infinitesimal.
A new \texttt{SurrealAlgebraicRoot} can use the same general idea of an
implicit exact object while supplying its own ordered-domain and branch
certificates. Native ordinary \texttt{Root} objects can still be used for
residual algebraic constants.

\subsection{\texttt{SeriesData} is a useful view, not the universal storage format}
A basic \texttt{SeriesData} object stores a finite coefficient list, integer
bounds, and a common denominator for the exponent progression. It carries
omitted-order information and supports arithmetic and calculus operations
\cite{WLSeriesData}. This is well suited to a finite rational-power jet.
It is not a representation of every well-ordered surreal exponent support.
The fact that \texttt{Series} can also produce certain logarithmic and other
expansions does not change that limitation of the basic exponent-lattice
format \cite{WLSeries}.

The extension should therefore export a compatible rank-one jet to
\texttt{SeriesData} when possible, and use its own support-aware object in
general. Converting with \texttt{Normal} discards the omitted-order term; it
must not be used internally as a proof that an infinite exact series equals
its displayed polynomial. Round-tripping a jet through plain syntax should
retain its precision certificate in a wrapper or explicitly announce its loss.

\subsection{Use the ordinary kernel where its assumptions actually apply}
Ordinary coefficient operations can be delegated directly when their inputs
are ordinary algebraic numbers or declared ordinary functions. Sparse
polynomial identities in formal monomial variables can also be computed by
the kernel, followed by the certified monomial interpretation. For ordered
questions over a surreal field, one route is to perform real-closed-field
quantifier elimination with symbolic parameters first and then evaluate the
resulting finitely many sign conditions using the custom ordered-field backend.

What should not be done is to ask native \texttt{Reduce} over the ordinary
reals to find a positive real $t$ smaller than $1/n$ for every positive
standard integer $n$. Such a real does not exist. Replacing that condition by
finitely many inequalities only specifies a small ordinary parameter, not
an infinitesimal. The mathematical transfer principle concerns finite
first-order field formulas in an appropriate real-closed model; it does not
turn unrestricted quantification over the integers into a field-language
assumption \cite{SourcePolynomial,PerrucciRoy}.

\subsection{Custom heads and conservative operator integration}
A first package can expose explicit operations such as
\texttt{SAdd}, \texttt{SMultiply}, \texttt{SInverse}, and
\texttt{SCompare}. This is easier to audit than widespread automatic
overloading. Once the core is stable, carefully scoped upvalues can offer
natural arithmetic notation. Wolfram documents upvalues as definitions
associated with arguments of an expression, and \texttt{Inactive} as a way
to suppress ordinary automatic evaluation \cite{WLUpValues,WLInactive}.

Those mechanisms do not create mathematical correctness on their own.
The package must avoid evaluation loops, ambiguous mixed-domain coercions,
and branch-changing rewrites. It should never redefine native
\texttt{Infinity} to mean $\omega$: the latter is a field element with
$\omega-\omega=0$, whereas the former has extended-limit semantics.

The numerical interface also needs restraint. \texttt{NumericQ} is connected
to ordinary numerical evaluation behavior \cite{WLNumericQ}. Making every
surreal wrapper numeric without defining a compatible numerical protocol can
send ordinary plotting or solving code into a false expectation of machine
or arbitrary-precision complex values. Separate scale-aware approximation
functions are preferable initially. \texttt{PossibleZeroQ}, documented as a
heuristic facility, must not be used as the proof of a leading coefficient's
nonvanishing or vanishing \cite{WLPossibleZeroQ}.

\section{A proposed software architecture and public contract}
\label{sec:architecture}

\subsection{Keep the domain, exact value, and computed view separate}
The recommended architecture has three different kinds of objects. A
\emph{domain} identifies the coefficient field, ordered exponent group,
embedding into the surreal numbers, and any exponential or derivation
conventions. An \emph{exact value} is a finite expression or certified
construction in that domain. A \emph{view} is a computed expansion, algebraic
root approximation, or enclosure of an exact value. Computing another view
must not silently change the value or its domain.

For example, a rational expression, an algebraic root description, and a
Hahn expansion can denote the same element while serving different purposes.
The rational expression is best for exact cancellation; the root description
is best for finite algebraic closure; the Hahn view is best for scale
comparisons. Keeping all three as related representations is more useful than
forcing every result into one preferred infinite normal form.

\begin{center}
\begin{tabularx}{\textwidth}{@{}P{0.25\textwidth}Y@{}}
\toprule
Component & Responsibilities and explicit guarantees\\
\midrule
Coefficient backend & Exact arithmetic; zero and sign decisions where
available; ordinary analytic identities and coefficient enclosures.\\[0.4em]
Exponent backend & Addition, rational division when supported, mathematical
order, equality, and order-preserving changes of scale.\\[0.4em]
Exact algebra & Sparse sums, rational expressions, finite algebraic
extensions, polynomial and matrix operations.\\[0.4em]
Support engine & Certified constructors, finite decomposition bounds,
coefficient production, and separate leading-term and prefix services.\\[0.4em]
Precision engine & Valuation bounds, generator-index jets, ordinary
coefficient error, refinement, and dependence tracking.\\[0.4em]
Function layer & Formal germs, ordinary domains, Hahn-coherent data,
branch selectors, and explicitly selected transserial interpretations.\\[0.4em]
Presentation layer & Readable notation, lossless serialization, evidence
inspection, and safe adapters to native CAS facilities.\\
\bottomrule
\end{tabularx}
\end{center}

\subsection{Capabilities should be queryable, not guessed from a head}
A domain should advertise separately whether equality and order are
decidable, whether coefficients and leading terms are computable, and
whether requested cutoff expansions are finite. Some capabilities
are local to a value or request: a particular series can have a certified
leading term even when its entire representation class has no general
leading-term algorithm. A cutoff may admit a finite expansion although a
larger cutoff does not.

The following is \emph{proposed interface syntax}, not functionality provided
by the demonstration package:
\begin{lstlisting}
dom = SurrealDomain[
  "Coefficients" -> "RealAlgebraic",
  "Exponents" -> LexicographicRationals[2],
  "Embedding" -> ConwayMonomials[{Omega, 1}]
];
x = ExactSurreal[representation, dom];
SurrealEqual[x, y, "Evidence" -> True]
SurrealExpand[x, "ValuationCutoff" -> {1, 0}]
SurrealCoefficient[x, {0, 17}]
SurrealRoot[p, "Selector" -> certifiedSelector]
CoherentLift[f, "Domain" -> U, "Branch" -> branchData]
\end{lstlisting}
The embedding record in this example means that exponent $(a,b)$ denotes
$a\omega+b$. It is not an instruction to identify an unspecified host
symbol with the entire surreal class. An implementation should store such
records in a validated internal format rather than rely on their printed
appearance.

A successful equality request can return a Boolean and a checkable witness.
An unsupported or exhausted request should return a distinct object such as
\texttt{Undetermined[reason, obligations]}. Invalid input should return a
failure. These are different outcomes: an unproved equality is not a proved
inequality, and a resource limit is not a mathematical counterexample.

\subsection{Certificates and immutable context}
Useful certificates include an exact coefficient normal form, a nonzero
leading term with a lower-support exclusion, an algebraic root selector, a
positive-grid decomposition bound, a nonvanishing ordinary denominator on a
specified domain, or a polynomial residual with a valuation lower bound.
A certificate checker can often be simpler than the search that found the
certificate. The checker need only verify the relevant finite statement;
it need not repeat an expensive simplification or root search.

Cache keys must include the coefficient assumptions, exponent order,
embedding, branch convention, and derivative interpretation. A normal form
computed under one convention must not be reused after any of those change.
In particular, the source's different $E_\alpha$ extensions cannot share a
cache keyed only by the printed expression \texttt{Exp[z]}.

Serialization should preserve these records and reference shared expression
nodes rather than copy exponentially growing trees. It should also preserve
whether a displayed term is exact or carries a remainder. Importing an
arbitrary coefficient program should not execute it automatically: mathematical
admissibility and safe program execution are separate checks. A restricted
constructor language, explicit resource budgets, and opt-in evaluation of
external code are suitable engineering boundaries.

\subsection{Lossless coercion before convenient notation}
Embedding a smaller ordered monomial field into a larger one is legitimate
when its map is explicit and order preserving. Adjoining algebraic roots is
also legitimate. Replacing a named transcendental constant by an independent
symbol, replacing an exact value by a jet, changing a branch, or discarding
an infinitesimal by taking standard part is not a lossless coercion.

The public interface should make those operations visibly intentional.
Taking a standard part, for example, requires a certificate of finiteness.
Scale specialization must distinguish ordinary parameter specialization from
evaluation at an actual surreal monomial. A finite algebraic identity
survives any defined field homomorphism, but an arbitrary infinite Hahn
expansion need not survive replacement of its scale by a small positive real
number. The nondiscrete-support and factorial-series examples already give
two independent reasons.

\section{An executed Wolfram Language prototype}
\label{sec:prototype}

\subsection{Exactly what the supplied code implements}
The archive contains \texttt{SurrealCASCore.wl}, a small package for finite
Hahn sums with Gaussian-rational coefficients and exponent vectors in
$\Q^r$ with lexicographic order. The rank is a positive ordinary integer.
The package uses explicit functions, validates its inputs, combines coincident
exponents, removes zero coefficients, and orders the remaining finite support
by a mathematical lexicographic comparator.

It implements addition, negation, multiplication, nonnegative integer powers,
conjugation, valuation, leading coefficient, and comparison of real-valued
inputs. A second representation stores a quotient of two such sums. It
implements fraction addition, multiplication, inversion, equality by cross
multiplication, and valuation. The denominator must be nonzero.

Thus the package handles a concrete fraction field of a finite-support group
algebra. Its exact fractions can denote infinitely supported Hahn elements
without generating those supports. The package does \emph{not} implement
algebraic coefficient fields, algebraic root selectors, infinite series
generators, certified analytic domains, arbitrary Conway cuts, or the
transserial layer. Those are architectural extensions discussed in this
article, not hidden functions in the delivered code.

\subsection{Sparse arithmetic and an infinite tail kept implicit}
A typical session begins as follows:
\begin{lstlisting}
Get["SurrealCASCore.wl"];

one = HCreate[{{{0}, 1}}, 1];
t   = HCreate[{{{1}, 1}}, 1];
a   = HAdd[one, HNegate[t]];
b   = HAdd[one, t];

HMultiply[a, b]
(* HData[1, {{{0}, 1}, {{2}, -1}}] *)

REqual[
  RCreate[one, a],
  RCreate[b, HMultiply[a, b]]
]
(* True *)
\end{lstlisting}
The first result is the exact equality $(1-t)(1+t)=1-t^2$.
The second checks
\[
 \frac1{1-t}=\frac{1+t}{1-t^2}
\]
without expanding either geometric series. No infinite summation routine is
needed for this equality. Fractions are not reduced by a polynomial gcd in
this demonstration, so structurally different fraction objects may be equal;
\texttt{REqual}, not structural equality of the wrappers, is the intended
comparison operation.

\subsection{Two genuinely different infinitesimal scales}
For rank two, interpret $(a,b)$ as $a\omega+b$ and use:
\begin{lstlisting}
s = HCreate[{{{0, 1}, 1}}, 2];
u = HCreate[{{{1, 0}, 1}}, 2];

HCompareReal[HPower[s, 100], u]
(* 1 *)
\end{lstlisting}
Here $s=\omega^{-1}$ and $u=\omega^{-\omega}$. The output means
$s^{100}>u$. More generally $s^n>u$ for every positive ordinary integer $n$,
because $(0,n)<(1,0)$ lexicographically. That general assertion follows from
the domain's order definition, not from extrapolating the particular test
with exponent $100$.

The constructor rejects machine-real coefficients such as $0.5$, irrational
exponent coordinates outside this backend, and inconsistent vector lengths.
Real comparison rejects nonreal coefficients. Inversion of zero and a zero
denominator are explicit failures. These rejection paths are part of the
tests, because a dependable exact core needs to protect its domain as well
as compute identities within it.

\subsection{Executed tests and what they establish}
The actual package file and its test harness were written into a temporary
directory and loaded in a Wolfram Language kernel. The returned report was:
\begin{lstlisting}
<|"Kernel" ->
    "15.0.1 for Linux x86 (64-bit) (July 2, 2026)",
  "Checks" -> 110, "Passed" -> 110, "Failures" -> {}|>
\end{lstlisting}
The tests cover deterministic arithmetic examples, fraction operations,
validation failures, and thirty reproducibly generated triples of sparse
Gaussian-rational sums. The generated cases check distributivity,
compatibility of conjugation with multiplication, and additivity of the
valuation on nonzero products. The archive includes the complete harness
\texttt{TestCore.wls} and a report describing the execution method.

A separate exact-symbolic script, \texttt{verify\_examples.py}, ran with
Python 3.13.5 and SymPy 1.14.0. All 69 checks passed. These verify the displayed
Newton and residual examples, root-expansion residual orders, finite residue
and multiplication matrices, several coefficient-extraction cases with
coincident grid exponents, and finite instances of the precision estimates.
The generated report records the individual checks.

These are tests of the delivered implementation and finite examples, not
proof-assistant verification of the general theorems. In particular, the
prototype does not execute real-closed-field quantifier elimination,
construct all algebraic roots, solve arbitrary support-recursive equations,
or certify the supplied manuscripts' analytic results. Its small scope makes
the tested functionality identifiable and reproducible.

\section{Implementation priorities, costs, and regression boundaries}
\label{sec:roadmap}

\subsection{Build in stages with useful stopping points}
An initial production-quality stage should improve the sparse and rational
core: algebraic coefficients, efficient multivariate rational normalization,
exponent-domain embeddings, expression sharing, and exact matrix and polynomial
adapters. This stage already supports many examples from the polynomial
manuscript, including resultants, discriminants, finite quotient algebras,
and initial polynomials for coefficients whose valuations are available.

The next stage adds real-algebraic root descriptions and their surcomplex
pairs. Its central deliverable is exact algebraic closure with robust branch
selection, not a universal root expansion. A subsequent series stage adds
rank-one finite truncations, rational-grid coefficient extraction, support
certificates, and precision propagation. High-rank requests should remain
compressed or return an informative unsupported-finite-view result when
necessary.

The analytic stage should begin with ordinary holomorphic lifts, rational
functions, regular implicit germs, and Hahn-coherent positive perturbations
with executable coefficient rules. It can then add preparation, division,
and residue methods from the supplied sources under their stated hypotheses.
Finally, an explicitly selected exp--log or transseries language can provide
additional scales and differential operations. This ordering is not a claim
that later stages are easy; it prevents their difficulties from blocking a
valuable exact algebraic system.

A CAS other than Wolfram Language can follow the same architecture. The
mathematical requirements are an exact coefficient engine, programmable
expression types, an ordered exponent representation, and validated adapters
for its polynomial and series facilities. Portability should reside in those
contracts and tests rather than in assumptions about a host's automatic
simplifier.

\subsection{Complexity must include representation and output size}
Multiplying sparse supports of sizes $m$ and $n$ initially produces at most
$mn$ pairs before collection. The bit lengths of rational coordinates and
algebraic coefficients matter as well as the number of terms. Repeated
cross multiplication can cause fraction growth, so gcd reduction, common
subexpression sharing, and delayed expansion are important improvements.
The demonstration's straightforward power routine is intentionally not an
optimized exponentiation algorithm.

For grid coefficient extraction, the finite box in \cref{thm:grid} proves
termination but can be large. Integer linear-equation methods, dynamic
programming, and memoized convolution can improve it. They must preserve
exact treatment of coincident sums. Polynomial root descriptions can have
large algebraic degrees; quantifier elimination is a sound foundation, not
a guarantee of low computational cost. Specialized Newton and Hensel methods
can provide far better useful views without replacing the exact fallback.

No algorithm can print an infinite cutoff prefix as a finite explicit term
list. This is an output-representation limitation, not merely poor complexity.
Likewise, no time budget transforms the halting reductions into decision
procedures. A useful system reports whether a request is expensive, outside
its implemented domain, impossible in the requested output format, or
undetermined by its current evidence.

\subsection{Regression cases should protect mathematical distinctions}
A regression suite should include more than successful expansions. It should
reject $\sum_{n\ge1}t^{1/n}$ as an unordered-support Hahn sum, reject the
geometric expansion of $\sum t^nz^n$ at $z=t^{-2}$ while permitting the
separate rational continuation, and retain the difference between an exact
factorial Hahn series and an ordinary divergent series after real
specialization.

It should also prevent a high-rank cutoff from being handled by an endless
rank-one loop, preserve the selected extension at $i\omega$, distinguish
Conway monomials from generic exponential powers, and refuse to decide
coefficient equality from insufficient numerical precision. Algebraic root
tests should include infinitesimally adjacent roots that ordinary decimal
neighborhoods cannot separate. Quotient-algebra tests should include repeated
roots, so that a nondegenerate residue pairing is not confused with a
nondegenerate trace pairing.

For each successful approximate calculation, verify the defining residual
and the claimed error certificate whenever that check is available. A displayed
root expansion without a branch certificate and residual bound is only a
candidate; one with those certificates can be a dependable computed view of
an exact root object.

\section{Conclusions: the representable subsets and their operations}
\label{sec:conclusion}

\subsection{Numbers and functions with meaningful exact descriptions}
A substantial subset is available immediately: finite normal forms over an
effective coefficient field, rational expressions in ordered monomials, and
all algebraic extensions of a specified effective ordered monomial field.
The real-closed fields $F_r$ and their algebraically closed complexifications
$C_r$ make this answer precise. Every value in these fields has a finite
algebraic description, even when its Hahn normal form is infinite.

Beyond that core, exact descriptions can include certified recurrence-defined
series, positive-grid Hahn generators, regular implicit solutions, ordinary
analytic lifts, Hahn-coherent functions on specified ordinary domains, and
selected transserial expressions. Their exactness means that their semantics
are specified; it does not imply that all their identities are decidable or
that every requested expansion is finite. Arbitrary formal ordinary-coefficient
series can also define exact values at infinitesimals when their coefficients
are supplied effectively, but unrestricted coefficient programs lose total
zero testing even in this simple setting.

There is no largest finite-description class that simultaneously solves all
these problems. Enlarging the expression language usually changes its
closure, equality, order, and evaluation guarantees. The right product is a
family of compatible exact domains with explicit capabilities.

\subsection{Operations that should be implemented, and their limits}
Exact field arithmetic, real comparison, conjugation, algebraic root
specification, finite polynomial algebra, and finite linear algebra are
realistic terminating operations in the effective algebraic core. Leading
terms, standard parts, Newton profiles, root-cluster counts, and certified
series refinement are realistic when the representation supplies the required
valuation and support services. Finite quotient algebras can often compute
residues, traces, and multiplicities without expanding individual roots.

Formal and coherent differentiation, regular implicit solving, positive-support
Taylor substitution, local inverses, preparation, and coefficientwise contour
operations can be implemented for suitably described functions. Global real
exponentiation can be admitted with its established surreal semantics;
global surcomplex transcendental operations must additionally state their
extension and branch policies. Integration must name its derivative and
function category.

The central design principle is therefore not to equate an infinite object
with an approximation, or a mathematical existence theorem with a program.
It is to preserve exact objects, expose certified finite information about
them, and make every change of representation, scale, domain, or analytic
interpretation explicit. That principle turns the supplied surcomplex theory
into a credible computer-algebra program rather than a collection of unsafe
rewrite rules.

\appendix
\section{Reproducibility and archive contents}
\label{app:repro}
The article is standalone LaTeX, with its bibliography included below. It uses
no external figures, data files, or bibliography database. Compile with:
\begin{lstlisting}
latexmk -pdf surreal_surcomplex_cas.tex
\end{lstlisting}
Alternatively run \texttt{pdflatex} until the references and table of contents
stabilize. The supplied PDF was compiled from the included source.

Run the Wolfram Language tests with:
\begin{lstlisting}
wolframscript -file TestCore.wls
\end{lstlisting}
The harness locates \texttt{SurrealCASCore.wl} relative to its own path, so the
working directory need not be the package directory. The recorded executed
kernel version is specified in \cref{sec:prototype}; compatibility with other
versions has not been separately tested.

Run the finite mathematical checks with:
\begin{lstlisting}
python verify_examples.py
\end{lstlisting}
The script requires SymPy and uses exact symbolic arithmetic, not floating-point
fits. It regenerates \texttt{python\_verification\_report.txt} beside the script.
The Wolfram report is a record of an actual remote kernel run, with its
provenance explained in \texttt{wolfram\_test\_report.txt}.

The archive contains the article source and PDF, the Wolfram package and test
harness, the Python verification script, both test reports, and a README.
The three input manuscripts are cited rather than bundled as dependencies.
Their general theorems have not been independently proof-assistant checked in
preparing this archive.

\section{Notation and semantic conventions}
\begingroup\small
\begin{longtable}{@{}P{0.25\textwidth}P{0.68\textwidth}@{}}
\toprule
Notation & Meaning\\
\midrule
\endhead
$\No,\SC$ & Surreal class field and its complexification $\No[i]$.\\[0.3em]
$t^\gamma$ & Conway monomial $\omega^{-\gamma}$, not a generic power operation.\\[0.3em]
$\RA,\CA$ & Real algebraic numbers and complex algebraic numbers.\\[0.3em]
$\Q^r_{\rm lex}$ & Rational exponent vectors, ordered by their first differing coordinate.\\[0.3em]
$E_r,F_r,C_r$ & Ordered rational monomial field, its relative real closure,
and the latter's complexification.\\[0.3em]
$v(A),\lc(A)$ & Least Hahn exponent and its coefficient; $v(0)=+\infty$.\\[0.3em]
$|z|$ & Surreal modulus $\sqrt{z\bar z}$, distinct from the valuation.\\[0.3em]
$\mm,\st$ & Infinitesimals and standard part on finite elements.\\[0.3em]
$E^*$ & Monoid of finite sums from a positive well-ordered exponent set $E$.\\[0.3em]
$\Oh\beta$ & A remainder with valuation at least $\beta$, including zero.\\[0.3em]
$U\sh$ & Finite surcomplex points with ordinary standard part in $U$.\\[0.3em]
$f\sh,\HH(U)$ & Ordinary analytic Taylor lift and Hahn-coherent coefficient data.\\[0.3em]
$D_z,D_T,\partial$ & Coordinate derivative, formal scale-variable derivative,
and a separately selected surreal-field derivation.\\[0.3em]
$E_\alpha$ & The source manuscript's explicitly phase-selected global
surcomplex exponential; not an unqualified canonical complex exponential.\\[0.3em]
$I_{a,\rho}(P)$ & Initial polynomial of a finite polynomial at a center and valuation scale.\\[0.3em]
$\lambda_P(H)$ & Top coefficient of the remainder of $H$ modulo monic $P$.\\
\bottomrule
\end{longtable}
\endgroup

\clearpage
\begin{thebibliography}{99}
\addcontentsline{toc}{section}{References}
\small

\bibitem{SourcePolynomial}
\emph{Surcomplex Polynomial Algebra: Order Geometry, Newton Profiles, and
Scale-Resolved Stability}. Supplied research manuscript dated September 21,
2026; file \texttt{surcomplex\_polynomial\_algebra(2).tex}. The workspaces,
initial-polynomial formulas, perturbation bounds, and finite residue algebra
are used with their stated hypotheses. No independent publication or proof
verification is inferred from the filename.

\bibitem{SourceResearch}
\emph{Hartogs Coherence, Finite Maps, and Residue Duality over the Surcomplex
Numbers}. Supplied research manuscript dated September 21, 2026; file
\texttt{surcomplex\_research.tex}. In particular, the positive-operator
calculus, fixed-domain division, coordinate-power perturbations, and finite
residue pairings.

\bibitem{SourceAnalysis}
\emph{Surcomplex Analysis: Infinitesimal Calculus, Hahn-Coherent Holomorphy,
and Contour Theory}. Supplied research manuscript dated September 21, 2026;
file \texttt{surcomplex\_analysis(3).tex}. Its separate notions of local germs,
coherent data, coefficientwise contours, and all-scale functions are retained.

\bibitem{Gonshor}
Harry Gonshor, \emph{An Introduction to the Theory of Surreal Numbers},
London Mathematical Society Lecture Note Series \textbf{110}, Cambridge
University Press, 1986.

\bibitem{Poonen}
Bjorn Poonen, ``Maximally complete fields,'' \emph{L'Enseignement
Math\'ematique} (2) \textbf{39} (1993), 87--106.
\href{https://math.mit.edu/~poonen/papers/amsval.pdf}{Author-hosted text}.
Corollary 4 provides the algebraic-closure input for divisible Hahn workspaces.

\bibitem{vdDE}
Lou van den Dries and Philip Ehrlich, ``Fields of surreal numbers and
exponentiation,'' \emph{Fundamenta Mathematicae} \textbf{167} (2001), 173--188.
\href{https://doi.org/10.4064/fm167-2-3}{doi:10.4064/fm167-2-3}.
Erratum, \textbf{168} (2001), 295--297. We use the established field and
exponential framework, not ordinal-length estimates affected by the erratum.

\bibitem{PerrucciRoy}
Daniel Perrucci and Marie-Fran\c coise Roy, ``Elementary recursive quantifier
elimination based on Thom encoding and sign determination,''
\emph{Annals of Pure and Applied Logic} \textbf{168} (2017), no.~8, 1588--1604.
\href{https://doi.org/10.1016/j.apal.2017.03.001}{doi:10.1016/j.apal.2017.03.001};
\href{https://arxiv.org/abs/1609.02879}{arXiv:1609.02879}.

\bibitem{Richardson}
Daniel Richardson, ``Some undecidable problems involving elementary functions
of a real variable,'' \emph{The Journal of Symbolic Logic} \textbf{33},
no.~4, 514--520.
\href{https://doi.org/10.2307/2271358}{doi:10.2307/2271358}.
The restrictions of the particular undecidability theorem matter; it is not
used to assert undecidability of every restricted exponential language.

\bibitem{vdH}
Joris van der Hoeven, \emph{Transseries and Real Differential Algebra},
Lecture Notes in Mathematics \textbf{1888}, Springer, 2006.
\href{https://www.texmacs.org/joris/ln/ln.html}{Author-hosted full text}.
The grid-based and support-controlled viewpoint informs the architecture;
the present demonstration does not implement that theory in full.

\bibitem{BMDerivations}
Alessandro Berarducci and Vincenzo Mantova, ``Surreal numbers, derivations and
transseries,'' \emph{Journal of the European Mathematical Society}
\textbf{20} (2018), 339--390.
\href{https://arxiv.org/abs/1503.00315}{arXiv:1503.00315}.

\bibitem{BMFunctions}
Alessandro Berarducci and Vincenzo Mantova, ``Transseries as germs of surreal
functions,'' \emph{Transactions of the American Mathematical Society}
\textbf{371} (2019), 3549--3592.
\href{https://doi.org/10.1090/tran/7428}{doi:10.1090/tran/7428};
\href{https://arxiv.org/abs/1703.01995}{arXiv:1703.01995}.

\bibitem{CostinEhrlich}
Ovidiu Costin and Philip Ehrlich, ``Integration on the surreals,''
\emph{Advances in Mathematics} \textbf{452} (2024), article 109823.
\href{https://arxiv.org/abs/2208.14331}{arXiv:2208.14331}, revised July 4, 2024.
This supplies a published extension and integration theory with its own
specified function classes, not an automatic interpretation of every Hahn sum.

\bibitem{WLRoot}
Wolfram Research, \emph{Root}, Wolfram Language documentation.
\href{https://reference.wolfram.com/language/ref/Root.html}{Official reference}.
Consulted September 21, 2026. Includes polynomial indexing and the documented
neighborhood-based root specifications.

\bibitem{WLAlgebraic}
Wolfram Research, \emph{AlgebraicNumber}, Wolfram Language documentation.
\href{https://reference.wolfram.com/language/ref/AlgebraicNumber.html}{Official reference}.
See also \href{https://reference.wolfram.com/language/tutorial/AlgebraicNumberFields.html}{Algebraic Number Fields}.
Consulted September 21, 2026.

\bibitem{WLRootReduce}
Wolfram Research, \emph{RootReduce}, Wolfram Language documentation.
\href{https://reference.wolfram.com/language/ref/RootReduce.html}{Official reference}.
Consulted September 21, 2026.

\bibitem{WLSeriesData}
Wolfram Research, \emph{SeriesData}, Wolfram Language documentation.
\href{https://reference.wolfram.com/language/ref/SeriesData.html}{Official reference}.
Consulted September 21, 2026.

\bibitem{WLSeriesCoefficient}
Wolfram Research, \emph{SeriesCoefficient}, Wolfram Language documentation.
\href{https://reference.wolfram.com/language/ref/SeriesCoefficient.html}{Official reference}.
Consulted September 21, 2026.

\bibitem{WLSeries}
Wolfram Research, \emph{Series}, Wolfram Language documentation.
\href{https://reference.wolfram.com/language/ref/Series.html}{Official reference}.
Consulted September 21, 2026.

\bibitem{WLInverseSeries}
Wolfram Research, \emph{InverseSeries}, Wolfram Language documentation.
\href{https://reference.wolfram.com/language/ref/InverseSeries.html}{Official reference}.
Consulted September 21, 2026.

\bibitem{WLAsymptotic}
Wolfram Research, \emph{Asymptotic}, Wolfram Language documentation.
\href{https://reference.wolfram.com/language/ref/Asymptotic.html}{Official reference}.
Consulted September 21, 2026.

\bibitem{WLAsymptoticSolve}
Wolfram Research, \emph{AsymptoticSolve}, Wolfram Language documentation.
\href{https://reference.wolfram.com/language/ref/AsymptoticSolve.html}{Official reference}.
Consulted September 21, 2026.

\bibitem{WLDifferentialRoot}
Wolfram Research, \emph{DifferentialRoot}, Wolfram Language documentation.
\href{https://reference.wolfram.com/language/ref/DifferentialRoot.html}{Official reference}.
Consulted September 21, 2026.

\bibitem{WLUpValues}
Wolfram Research, \emph{UpValues}, Wolfram Language documentation.
\href{https://reference.wolfram.com/language/ref/UpValues.html}{Official reference}.
Consulted September 21, 2026.

\bibitem{WLInactive}
Wolfram Research, \emph{Inactive}, Wolfram Language documentation.
\href{https://reference.wolfram.com/language/ref/Inactive.html}{Official reference}.
Consulted September 21, 2026.

\bibitem{WLNumericQ}
Wolfram Research, \emph{NumericQ}, Wolfram Language documentation.
\href{https://reference.wolfram.com/language/ref/NumericQ.html}{Official reference}.
Consulted September 21, 2026.

\bibitem{WLPossibleZeroQ}
Wolfram Research, \emph{PossibleZeroQ}, Wolfram Language documentation.
\href{https://reference.wolfram.com/language/ref/PossibleZeroQ.html}{Official reference}.
Consulted September 21, 2026.
\end{thebibliography}
\end{document}
